\documentclass[12pt]{article}
\usepackage[T1]{fontenc}
\usepackage[margin=1in]{geometry}
\usepackage{amsmath,amssymb,amsthm,mathtools}
\usepackage{booktabs,longtable,array}
\usepackage{setspace}
\usepackage{xurl}
\usepackage{hyperref}
\usepackage{enumitem}

\setstretch{1.15}
\setlength{\parskip}{0.4em}
\setlength{\parindent}{0pt}

\newtheorem{assumption}{Assumption}
\newtheorem{definition}{Definition}
\newtheorem{lemma}{Lemma}
\newtheorem{proposition}{Proposition}
\newtheorem{corollary}{Corollary}
\theoremstyle{remark}
\newtheorem{remark}{Remark}

\newcommand{\E}{\mathbb E}
\newcommand{\R}{\mathbb R}
\newcommand{\dd}{\,d}
\newcommand{\mc}{\mathrm{mc}}
\newcommand{\M}{\mathcal M}
\newcommand{\Q}{\mathcal Q}
\newcommand{\V}{\mathcal V}
\newcommand{\C}{\mathcal C}
\newcommand{\one}{\mathbf 1}

\title{Baseline Model: Monetary Policy, Critical Imported Inputs, and Endogenous Repair}
\author{}
\date{}

\begin{document}

\maketitle

\section{Purpose}

This note develops the baseline model for an economy exposed to a persistent disruption in access
to a critical imported input. The model is deliberately organized around one mechanism. An external
bottleneck state makes the imported input more expensive to procure and less physically available;
when the availability restriction binds, the economy assigns an endogenous scarcity rent to the
limited input stock. Firms can reduce future exposure to that bottleneck by investing in
adaptation, so recovery is not a mechanical return to normal conditions but a dynamic production
choice. The baseline isolates one bottleneck-repair mechanism and studies its implications for
current scarcity, inflation dynamics, and firms' incentives to reduce future dependence on the
constrained input.

This construction differs from a reduced-form cost-push specification because the imported-input
distortion is not inserted directly as an exogenous markup wedge in firms' marginal cost. The
distortion is generated by a physical input constraint, its equilibrium shadow price, and firms'
chosen exposure to that constraint. The rest of the model builds these objects in sequence: the
external environment determines input conditions, households provide demand and discounting, firms
choose production and prices, and the adaptation problem links today's policy stance to tomorrow's
input dependence.

The intended economic environment is one in which a previously routine imported input becomes a
bottleneck. The input can be a component, software license, machine part, payment service,
logistics service, or imported technology without which domestic production becomes more costly or
less scalable. The model therefore treats the shock not as ``higher costs'' in the abstract, but as
a disruption in access to a concrete production requirement.

The model is positioned between three strands of work. The monetary-policy comparison follows the
logic of persistent supply-regime models such as Galo--Nu\~no, Renner, and Scheidegger, where the
natural benchmark becomes state dependent. The production side is closer to the bottleneck and
imported-input literature associated with nonlinear production networks and low substitutability of
foreign inputs. The new object added here is the repair decision: firms do not only face a bad
state; they can invest in reducing future dependence on the constrained input. This is why the
baseline separates the external procurement price $p_t^m$, the endogenous scarcity rent $\chi_t$,
and the adaptation stock $A_t$ from the beginning.

\part{Baseline Economy}

\section{Environment and Timing}

Time is discrete, $t=0,1,2,\ldots$. The economy is populated by a representative household, a
continuum of monopolistically competitive goods producers indexed by $i\in[0,1]$, a competitive
market for a critical imported input, and a monetary authority. A competitive CES aggregator
bundles differentiated varieties into the final consumption good. The aggregator is a demand system,
not a separate economic decision maker.

At the beginning of period $t$, the external disruption and relief stocks $(D_t,X_t)$ are observed.
These states summarize sanctions, logistics disruptions, payment restrictions, insurance and
compliance costs, workaround routes, settlement channels, and other external conditions governing
access to the critical imported input. The within-period timing is:
\begin{enumerate}[leftmargin=1.5em]
    \item the aggregate external, productivity, competitive supply-price, installed adaptation,
    and inherited price-dispersion states are observed;
    \item the policy instrument $R_t$ is set by the monetary-policy environment specified after the
    private baseline economy is defined;
    \item households choose consumption, labor supply, and nominal bond holdings;
    \item firms that receive a Calvo adjustment opportunity choose reset prices, while all firms
    choose static input demands and production taking the scarcity rent $\chi_t$ as a parametric
    market price;
    \item the imported-input market clears subject to the aggregate cap $\bar m_t$, and equilibrium
    selects the value of $\chi_t$ that is consistent with those cost-minimizing demands;
    \item firms choose adaptation investment $I_t^A$, which determines $A_{t+1}$;
    \item profits, scarcity-rent rebates, financing payments, and transfers are paid to the
    representative household.
\end{enumerate}
This timing makes the policy experiment below a current-period policy innovation in a given
external bottleneck environment. The external state determines input conditions before private
decisions are made, while adaptation chosen in period $t$ affects imported-input dependence from
period $t+1$ onward.

Equivalently, the period starts with the outside procurement environment. The stocks $(D_t,X_t)$
determine the external procurement price $p_t^m$ and the physical availability cap $\bar m_t$.
Given these conditions, firms treat $\chi_t$ parametrically as the price of current scarce access,
just as they take other aggregate prices as given; equilibrium then selects the value of $\chi_t$
that clears the imported-input market. Static input choices, marginal costs, and Calvo pricing are
determined at the current technology $A_t$, while adaptation investment is chosen
intertemporally because it changes the firm's future dependence on the bottlenecked input.
Institutionally, $\chi_t$ can be read as the market-clearing price of scarce access: the value of
an import permit, queue priority, delivery slot, or internal allocation right attached to a limited
input stock.

\begin{assumption}[External disruption and relief states]\label{ass:DX}
The external bottleneck environment is summarized by two stocks:
\begin{equation}
    D_t\ge 0,
    \qquad
    X_t\ge 0.
    \label{eq:DX_states}
\end{equation}
The stock $D_t$ measures accumulated external disruption pressure: sanctions, payment failures,
logistics breakdowns, customs restrictions, compliance costs, forced rerouting, and supplier exits.
It is a state of the outside procurement environment, not a firm-level productivity shock. The
stock $X_t$ measures accumulated external relief: new trade routes, restored settlement channels,
successful rerouting, partial easing, alternative suppliers, and adaptation of outside delivery
infrastructure. The net external bottleneck tightness is
\begin{equation}
    \mathcal B_t\equiv D_t-X_t.
    \label{eq:B_definition}
\end{equation}
\end{assumption}

The distinction between $D_t$ and $X_t$ is not cosmetic. Disruptions and relief do not have the
same timing or durability. Damage can arrive sharply, cluster, and decay slowly. Relief usually
arrives with delay, is partial, and may be more fragile. Firm adaptation $A_t(i)$ is kept separate
from both objects: $D_t$ and $X_t$ describe the outside procurement environment, while $A_t(i)$
describes the firm's own technology.

In applications, $D_t$ should be read as accumulated external damage rather than a one-period
shock. It is high when firms face more blocked routes, rejected payments, unavailable suppliers, or
compliance barriers. The relief stock $X_t$ captures the opposite side of the same external
environment: new intermediaries, alternative routes, partial licensing, or settlement channels.
These objects are external to the firm. They describe what the outside world allows the firm to do.

The next equations specify this external environment. They should be read as one stochastic block:
damage arrives through marked events and decays slowly, relief also arrives through marked events
but may be less durable, and event intensities can themselves depend on the accumulated state of
the disruption. The assumptions are separated only to keep the roles of damage, relief, and
arrival intensities transparent.

\begin{assumption}[Disruption stock]\label{ass:D_law}
The disruption stock evolves according to
\begin{equation}
    D_{t+1}
    =
    (1-\delta_D)D_t
    +
    U_{t+1}^D,
    \qquad
    \delta_D\in(0,1),
    \label{eq:D_law}
\end{equation}
where
\begin{equation}
    U_{t+1}^D
    =
    \sum_{j=1}^{N_{t+1}^D}m_{t+1,j}^D,
    \qquad
    N_{t+1}^D\mid\lambda_t^D
    \sim
    \mathrm{Poisson}(\lambda_t^D),
    \qquad
    m_{t+1,j}^D>0.
    \label{eq:D_marks}
\end{equation}
The decay rate $\delta_D$ is small in persistent disruption episodes.
\end{assumption}

The marked-Poisson specification captures the fact that adverse external restrictions often arrive
as events rather than smooth Gaussian innovations. A new sanction package, a payment-system
failure, a shipping closure, or a supplier exit can discretely raise the disruption stock. The
effect then decays slowly as firms and intermediaries learn to operate around it.

This timing is deliberately different from a standard AR(1) cost shock. In an AR(1) process, bad
conditions mechanically fade every period. Here, bad conditions can persist because old restrictions
decay slowly and new restrictions may arrive before the old ones disappear.

\begin{assumption}[Relief stock]\label{ass:X_law}
The external relief stock evolves according to
\begin{equation}
    X_{t+1}
    =
    (1-\delta_X)X_t
    +
    U_{t+1}^X,
    \qquad
    \delta_X\in(0,1),
    \label{eq:X_law}
\end{equation}
where
\begin{equation}
    U_{t+1}^X
    =
    \sum_{j=1}^{N_{t+1}^X}\zeta_{t+1,j}^X,
    \qquad
    N_{t+1}^X\mid\lambda_t^X
    \sim
    \mathrm{Poisson}(\lambda_t^X),
    \qquad
    \zeta_{t+1,j}^X>0.
    \label{eq:X_marks}
\end{equation}
The empirically relevant case allows $\delta_X\ge \delta_D$, so that relief is weakly less durable
than disruption.
\end{assumption}

Relief arrivals capture partial external normalization. They do not mean that the original supply
chain is fully restored. They may represent narrow workarounds: a new route, a new payment channel,
a new intermediary, or a foreign supplier that can deliver only at limited scale.

\begin{assumption}[Event intensities]\label{ass:intensities}
The event intensities may be constant, or they may evolve according to
\begin{align}
    \log\lambda_{t+1}^D
    &=
    (1-\rho_{\lambda D})\log\bar\lambda_D
    +
    \rho_{\lambda D}\log\lambda_t^D
    +
    \kappa_D^\lambda D_t
    +
    \sigma_{\lambda D}\varepsilon_{t+1}^{\lambda D},
    \label{eq:lambda_D_law}\\
    \log\lambda_{t+1}^X
    &=
    (1-\rho_{\lambda X})\log\bar\lambda_X
    +
    \rho_{\lambda X}\log\lambda_t^X
    +
    \beta_X D_t
    +
    \sigma_{\lambda X}\varepsilon_{t+1}^{\lambda X},
    \qquad
    \beta_X>0.
    \label{eq:lambda_X_law}
\end{align}
The coefficient $\kappa_D^\lambda\ge 0$ allows disruption to be self-reinforcing: a bad external environment
can raise the probability of further adverse events. The coefficient $\beta_X>0$ allows accumulated
disruption to trigger relief attempts, but only through the arrival intensity, not through an
instantaneous one-for-one offset.

The relief intensity does not mean that damage automatically repairs itself. It means that a more
severe disruption creates incentives for external workarounds to appear: traders search for new
routes, banks search for settlement channels, and suppliers reorganize delivery chains. These
responses arrive probabilistically and may still be too small to restore the pre-shock environment.
\end{assumption}

\begin{assumption}[Mapping from external states to input conditions]\label{ass:DX_mapping}
The external stocks determine the procurement price and the available quantity of the critical
imported input:
\begin{align}
    \log p_t^m
    &=
    \log\bar p^m
    +
    \nu_{pD}D_t
    -
    \nu_{pX}X_t,
    \label{eq:pm_DX}\\
    \log\bar m_t
    &=
    \log\bar m
    -
    \nu_{qD}D_t
    +
    \nu_{qX}X_t,
    \label{eq:mbar_DX}
\end{align}
where all $\nu$ coefficients are nonnegative.
\end{assumption}

Equations \eqref{eq:pm_DX}--\eqref{eq:mbar_DX} are the point at which the external environment
enters firms' real decisions. A higher disruption stock can make the input more expensive to buy
and also reduce the amount that can physically be obtained. A higher relief stock works in the
opposite direction. This distinction matters because a high price and a hard quantity ceiling have
different equilibrium implications.

The mapping allows asymmetric effects. The restriction
\begin{equation}
    \nu_{qD}\ge \nu_{pD}
    \label{eq:D_asymmetry}
\end{equation}
captures cases in which disruption hits physical availability especially strongly. The restriction
\begin{equation}
    \nu_{qX}>\nu_{pX}
    \label{eq:X_asymmetry}
\end{equation}
captures cases in which relief restores access more than it lowers procurement costs: new routes
may make delivery possible again without making it cheap.

For example, an alternative shipping route can restore physical delivery while remaining slow,
risky, and expensive. In that case $\bar m_t$ recovers more than $p_t^m$ falls. This is why the
model keeps the price and quantity channels separate instead of collapsing them into one wedge.

\begin{remark}[External relief versus firm adaptation]\label{rem:relief_vs_adaptation}
External relief $X_t$ improves the outside procurement environment and enters
\eqref{eq:pm_DX}--\eqref{eq:mbar_DX}. Firm adaptation $A_t(i)$ does not change $D_t$, $X_t$,
$p_t^m$, or $\bar m_t$. It changes the firm's production technology and reduces dependence on the
imported input. Thus relief and adaptation are distinct recovery forces: relief changes the
external environment, while adaptation changes domestic exposure to that environment.
\end{remark}

The recursive external state is therefore
\begin{equation}
    s_t^B=(D_t,X_t,\lambda_t^D,\lambda_t^X),
    \label{eq:external_state}
\end{equation}
with $\lambda_t^D$ and $\lambda_t^X$ omitted when intensities are constant. This state summarizes
the external side of the model: persistence comes from slowly decaying stocks, clustered adverse
events enter through the disruption arrivals, delayed relief enters through the relief arrivals,
and the mapping to $(p_t^m,\bar m_t)$ determines whether the same external episode mainly raises
procurement costs, tightens physical availability, or does both.

\begin{assumption}[Admissible stochastic environment]\label{ass:admissible_env}
The exogenous state vector, including $(D_t,X_t,\lambda_t^D,\lambda_t^X)$, productivity $Z_t$, and
the competitive supply prices $(p_t^d,p_{A,t})$, evolves on an admissible state space
$\mathcal S$ such that $p_t^m>0$, $\bar m_t>0$, $Z_t>0$, $p_t^d>0$, $p_{A,t}>0$, and all
conditional expectations below are finite. Productivity and competitive supply prices are positive
Markov processes with transition laws taken as given by private agents.
\end{assumption}

Assumption~\ref{ass:admissible_env} closes the external block by ruling out explosive paths on
which the procurement price, available input capacity, or productivity leaves the economically
meaningful domain. With the external state now defined, the next step is to specify how private
agents make decisions given those input conditions.

\section{Households}

The representative household has preferences
\begin{equation}
    \E_0\sum_{t=0}^\infty \beta^t
    \left[
        \frac{C_t^{1-\sigma}}{1-\sigma}
        -
        \frac{N_t^{1+\varphi}}{1+\varphi}
    \right],
    \qquad
    \beta\in(0,1),\quad \sigma>0,\quad \varphi>0.
    \label{eq:hh_preferences}
\end{equation}
The nominal budget constraint is
\begin{equation}
    P_t C_t + B_t^h
    \le
    W_t N_t + R_{t-1}B_{t-1}^h + \Pi_t^f + T_t,
    \label{eq:hh_budget_nominal}
\end{equation}
where $P_t$ is the final-good price index, $W_t$ is the nominal wage, $B_t^h$ denotes one-period
nominal bond holdings, $\Pi_t^f$ is nominal firm profit rebated to the household, and $T_t$ is a
lump-sum transfer. Let real wages be $w_t=W_t/P_t$, gross inflation be
\begin{equation}
    \Pi_{t+1}=\frac{P_{t+1}}{P_t},
    \label{eq:inflation_def}
\end{equation}
and the household's marginal utility of consumption be
\begin{equation}
    \Lambda_t=C_t^{-\sigma}.
    \label{eq:lambda_def}
\end{equation}

\begin{lemma}[Household optimality]\label{lem:hh_focs}
If the household chooses an interior allocation, then
\begin{align}
    C_t^{-\sigma} &= \Lambda_t, \label{eq:hh_muc}\\
    N_t^\varphi &= w_t\Lambda_t, \label{eq:hh_labor}\\
    1 &= \beta R_t\E_t
    \left[
        \frac{\Lambda_{t+1}}{\Lambda_t}
        \frac{1}{\Pi_{t+1}}
    \right].
    \label{eq:hh_euler}
\end{align}
The stochastic discount factor for real payoffs is
\begin{equation}
    M_{t,t+1}
    =
    \beta\frac{\Lambda_{t+1}}{\Lambda_t}.
    \label{eq:sdf}
\end{equation}
\end{lemma}

The household block provides both aggregate demand and the discounting kernel used in firm
valuation. The Euler equation is the channel through which the nominal policy rate affects current
spending: a higher $R_t$ raises the return to postponing consumption, all else equal. The same
stochastic discount factor $M_{t,t+1}$ is then used by firms when they value future cost savings
from adaptation, so the household side links monetary policy to both demand and repair incentives.

The external bottleneck does not enter household preferences as a separate utility shock. Households
are affected through equilibrium prices and quantities. When disruption raises the effective cost
of the critical input or makes the input scarce, firms' marginal costs, inflation, output, wages,
and profits adjust; these objects determine the household's consumption, labor supply, and real
income. Thus households bear the supply disruption through the allocation generated by the
bottleneck economy rather than through an additional taste or endowment disturbance.
The intertemporal-saving channel is still present through \eqref{eq:hh_euler}: expected future
consumption and inflation determine the real return households require to shift consumption across
periods. What changes relative to a pure reduced-form supply-regime model is the role of this
channel. Here it provides demand transmission and discounting for future repair payoffs, while the
new mechanism operates through imported-input scarcity and firms' incentives to reduce future
exposure to that scarcity.

\begin{proof}
The Lagrangian for \eqref{eq:hh_budget_nominal} has multiplier $\lambda_t^n$ on nominal resources.
The consumption first-order condition gives $C_t^{-\sigma}=\lambda_t^n P_t$. Define
$\Lambda_t=\lambda_t^nP_t$. The labor condition gives $N_t^\varphi=\lambda_t^n W_t=w_t\Lambda_t$.
The bond condition gives $\lambda_t^n=\beta R_t\E_t\lambda_{t+1}^n$. Multiplying by $P_t$ and using
$\lambda_{t+1}^n=\Lambda_{t+1}/P_{t+1}$ yields \eqref{eq:hh_euler}.
\end{proof}

\section{Final Demand and Price Index}

The final good is a CES aggregate of differentiated varieties:
\begin{equation}
    Y_t=
    \left[
        \int_0^1 y_t(i)^{\frac{\varepsilon-1}{\varepsilon}}\dd i
    \right]^{\frac{\varepsilon}{\varepsilon-1}},
    \qquad
    \varepsilon>1.
    \label{eq:final_aggregator}
\end{equation}
Cost minimization by the aggregator implies the demand curve
\begin{equation}
    y_t(i)
    =
    \left(\frac{P_t(i)}{P_t}\right)^{-\varepsilon}Y_t
    \label{eq:variety_demand}
\end{equation}
and the price index
\begin{equation}
    P_t=
    \left[
        \int_0^1 P_t(i)^{1-\varepsilon}\dd i
    \right]^{\frac{1}{1-\varepsilon}}.
    \label{eq:price_index}
\end{equation}

The CES aggregator is not a separate substantive agent. It is a compact way to obtain downward
sloping demand for each differentiated variety. This demand system gives individual firms
price-setting power; the next block therefore specifies how those firms produce, how the
critical-input bottleneck enters their costs, and how price rigidity transmits those costs into
inflation.

\section{Firms}

\subsection{Technology}

Given the demand curve above, each firm chooses inputs and prices taking aggregate conditions as
given. The production side is where the external bottleneck becomes economically concrete: firms
need a critical imported input, can use a costly substitute, and can gradually alter the technology
that governs the mix between them.

Each firm produces a differentiated good with labor $n_t(i)$ and an intermediate input composite
$x_t(i)$:
\begin{equation}
    y_t(i)=Z_t n_t(i)^{1-\alpha}x_t(i)^\alpha,
    \qquad
    \alpha\in(0,1),
    \label{eq:production}
\end{equation}
where $Z_t>0$ is aggregate productivity. The intermediate composite combines the critical imported
input $m_t(i)$ and a domestic or otherwise non-critical substitute input $d_t(i)$:
\begin{equation}
    x_t(i)=
    \left[
        \omega(A_t(i))m_t(i)^{\frac{\rho-1}{\rho}}
        +
        \bigl(1-\omega(A_t(i))\bigr)d_t(i)^{\frac{\rho-1}{\rho}}
    \right]^{\frac{\rho}{\rho-1}},
    \qquad
    \rho>0,\quad \rho\neq 1.
    \label{eq:x_ces}
\end{equation}
The parameter $\rho$ is the elasticity of substitution between the imported input and the
substitute input. Low $\rho$ means that the imported input is hard to replace in current
production. High $\rho$ means that substitution is technologically easier.

This elasticity is economically central. If $\rho$ is low, the firm cannot easily replace the
critical imported input with domestic or non-critical inputs in current production. Then the same
external bottleneck produces a larger rise in marginal cost and a stronger incentive to invest in
future adaptation. If $\rho$ is high, the firm can substitute away more easily even before new
adaptation is installed.

It is useful to distinguish this current substitution elasticity from dynamic adaptation. The
elasticity $\rho$ governs what the firm can do within the current technology when relative input
prices change. The stock $A_t(i)$ governs how that technology itself has been modified by past
repair spending. Thus a firm can respond to a bottleneck immediately by changing its current input
mix, but it reduces future exposure only by accumulating adaptation capital.

The stock $A_t(i)$ is firm-level adaptation capital. It shifts the intermediate-input technology
away from the imported input:
\begin{assumption}[Adaptation technology]\label{ass:omega}
The import-dependence weight $\omega(A)$ satisfies
\begin{equation}
    \omega(A)\in(0,1),\qquad
    \omega_A(A)<0,\qquad
    \omega_{AA}(A)\ge 0.
    \label{eq:omega_assumption}
\end{equation}
Thus additional adaptation reduces the technological weight placed on the critical imported input,
with weakly diminishing effectiveness.
\end{assumption}

The stock $A_t(i)$ should be interpreted as accumulated ability to use substitute inputs instead
of the critical imported input. In applications this includes product redesign, compatibility
adjustments, supplier switching, recertification, engineering changes, process reorganization,
localization of components, inventory systems, and workaround logistics. It does not increase the
external availability of the imported input. It reduces how much of that input the firm needs for a
given level of production, which is why it lowers $\omega(A_t(i))$ rather than changing
$(D_t,X_t)$. The firm is not simply buying less imported input in a given period; it is installing
the capability to operate with a different input mix in future production.

The substitute input $d_t(i)$ is available at real price $p_t^d>0$. In the baseline, $p_t^d$ is
taken as a competitive price process independent of the imported-input bottleneck. This isolates
the bottleneck mechanism: the model asks how firms respond when one critical input becomes
expensive or scarce, while the non-critical substitute remains available at a market price.
The substitute need not be the same good produced domestically. It is a broad class of imperfect
alternatives: domestic components, alternative materials, non-critical analogues, modified
production processes, or re-engineered inputs. The bottleneck remains economically relevant because
using these alternatives effectively requires the accumulated adaptation stock $A_t(i)$.

\begin{assumption}[Competitive supply prices]\label{ass:competitive_supply_prices}
The prices $p_t^d>0$ and $p_{A,t}>0$ are exogenous competitive supply prices for, respectively,
the non-critical substitute input and the adaptation installation good. They are Markov processes
on the admissible state space and are taken as given by households and firms.
\end{assumption}

Assumption~\ref{ass:competitive_supply_prices} is a baseline discipline, not a claim that domestic
substitute and repair resources are costless. For resource accounting, the firm pays the real
resource cost $p_t^d$ for each unit of the substitute input and $p_{A,t}$ for each unit of
adaptation installation. What is left outside the baseline is the production side of these inputs.
Therefore the baseline can speak about costly substitution and costly repair, but it does not yet
model domestic resource competition in producing substitute inputs or repair services. It therefore
does not generate wage pressure, factor reallocation, or repair-sector inflation from those
activities. Those effects belong to a later resource-based repair extension.

\subsection{Imported-Input Market and Scarcity Rent}

The critical imported input is available in aggregate amount $\bar m_t$ from
\eqref{eq:mbar_DX}. Individual firms take the effective real price of the imported input as given,
including the scarcity rent $\chi_t$. The aggregate imported-input feasibility constraint is
\begin{equation}
    \int_0^1 m_t(i)\dd i\le \bar m_t.
    \label{eq:m_cap}
\end{equation}
Let $\chi_t\ge 0$ denote the real scarcity rent associated with this constraint.
Equivalently, $\chi_t$ is the market-clearing price of scarce access to the imported input. Each
individual firm takes it as parametric, while the equilibrium value of $\chi_t$ is the one that
makes aggregate desired imported-input use consistent with the physical cap.

\begin{definition}[Effective imported-input price]\label{def:effective_m_price}
The effective real unit cost of the critical imported input is
\begin{equation}
    \tilde p_t^m=p_t^m+\chi_t.
    \label{eq:effective_m_price}
\end{equation}
The procurement component $p_t^m$ is the external price of obtaining the imported input from
outside the economy under the current disruption and relief states. The rent $\chi_t$ is the
domestic scarcity value that appears only when the limited available quantity must be allocated
across firms. Thus $\tilde p_t^m$ is not the sum of two arbitrary wedges: it is external
procurement cost plus the internal congestion value of scarce access.
\end{definition}

Market clearing requires the complementarity condition
\begin{equation}
    \chi_t\ge 0,\qquad
    \bar m_t-\int_0^1m_t(i)\dd i\ge 0,\qquad
    \chi_t\left(\bar m_t-\int_0^1m_t(i)\dd i\right)=0.
    \label{eq:m_complementarity}
\end{equation}

\begin{assumption}[Ownership and rebate of scarcity rents]\label{ass:rent_rebate}
Rights to the scarce imported-input capacity are ultimately owned by the representative household,
either directly, through firms, or through a public licensing authority. Scarcity-rent payments
$\chi_t\int_0^1m_t(i)\dd i$ are rebated lump-sum through firm profits or transfers.
\end{assumption}

When \eqref{eq:m_cap} is slack, $\chi_t=0$ and firms pay only the external procurement price
$p_t^m$. When \eqref{eq:m_cap} binds, $\chi_t>0$ raises the private cost of imported-input use until
aggregate demand for $m_t$ equals available supply.

The rent $\chi_t$ is therefore the model's measure of current scarcity. It is the value of one
additional unit of the constrained input inside the domestic economy. In applied terms, it can be
read as the shadow value of access rights, queue priority, import permits, scarce delivery slots,
or internal allocation power over a limited input. Assumption~\ref{ass:rent_rebate} determines how
that value is accounted for: it affects firms' private marginal costs and prices, but it is not a
resource cost for the representative household. Once this effective input price is determined, the
firm's static cost-minimization problem pins down how much labor, imported input, and substitute
input are needed for any chosen level of output.

\subsection{Static Cost Minimization}

This problem is static only in the sense that the firm takes the installed adaptation stock
$A_t(i)$ as predetermined within period $t$. It chooses the least-cost current input mix for a given
output level. The intertemporal choice of how much to change $A_t(i)$ is introduced separately in
the adaptation block below.

For fixed $A_t(i)$, $\tilde p_t^m$, $p_t^d$, and $w_t$, define the unit cost of the intermediate
composite $x_t(i)$ by
\begin{equation}
    p_t^x(A_t(i),\tilde p_t^m,p_t^d)
    =
    \left[
        \omega(A_t(i))^\rho(\tilde p_t^m)^{1-\rho}
        +
        \bigl(1-\omega(A_t(i))\bigr)^\rho(p_t^d)^{1-\rho}
    \right]^{\frac{1}{1-\rho}}.
    \label{eq:px_dual}
\end{equation}
This is the cheapest cost of producing one unit of the intermediate composite from the imported
input and the substitute input. It is the first place where current substitution possibilities
enter costs: when the imported input becomes more expensive through $p_t^m$ or $\chi_t$, the
firm can partially substitute toward $d_t(i)$, but the strength of that substitution is governed by
$\rho$ and by the installed adaptation stock through $\omega(A_t(i))$.
The real unit production cost of variety $i$ is
\begin{equation}
    \mc_t(i)
    =
    \frac{1}{Z_t}
    \left(\frac{w_t}{1-\alpha}\right)^{1-\alpha}
    \left(\frac{p_t^x(A_t(i),\tilde p_t^m,p_t^d)}{\alpha}\right)^\alpha.
    \label{eq:unit_cost}
\end{equation}
The marginal-cost object $\mc_t(i)$ is therefore not an exogenous markup wedge. It is the real cost
of producing one additional unit of variety $i$ after the firm has chosen the cheapest labor,
imported-input, and substitute-input mix at the current bottleneck conditions. This is the object
that later enters Calvo pricing and transmits the bottleneck into inflation.

\begin{lemma}[Cost-minimizing input demands]\label{lem:input_demands}
For a firm producing $y_t(i)$, static cost minimization at
$(A_t(i),\tilde p_t^m,p_t^d,w_t,Z_t)$ implies
\begin{align}
    n_t(i)
    &=
    (1-\alpha)\frac{\mc_t(i)y_t(i)}{w_t},
    \label{eq:n_demand}\\
    x_t(i)
    &=
    \alpha\frac{\mc_t(i)y_t(i)}{p_t^x(A_t(i),\tilde p_t^m,p_t^d)},
    \label{eq:x_demand}\\
    m_t(i)
    &=
    x_t(i)\,
    \omega(A_t(i))^\rho
    \left(
        \frac{p_t^x(A_t(i),\tilde p_t^m,p_t^d)}{\tilde p_t^m}
    \right)^\rho,
    \label{eq:m_demand}\\
    d_t(i)
    &=
    x_t(i)\,
    \bigl(1-\omega(A_t(i))\bigr)^\rho
    \left(
        \frac{p_t^x(A_t(i),\tilde p_t^m,p_t^d)}{p_t^d}
    \right)^\rho.
    \label{eq:d_demand}
\end{align}
\end{lemma}

The demand formulas show how the bottleneck enters firm costs. The external price $p_t^m$ and the
scarcity rent $\chi_t$ enter only through the effective imported-input price
$\tilde p_t^m$. When $\chi_t$ rises, firms behave as if the imported input became more expensive
even if the foreign procurement price did not change. Adaptation lowers the import-dependence
weight and therefore reduces exposure to this effective price.

These demands also separate two channels that are easy to conflate. A fall in output reduces
imported-input use mechanically through the scale of production $y_t(i)$; this is the current
bottleneck-relief channel used by monetary tightening. A rise in installed adaptation reduces
imported-input use through the intensity term multiplying $x_t(i)$; this is structural repair.
The first channel uses less of the constrained input because the firm produces less today. The
second channel uses less of the constrained input because the firm has changed how production is
organized.

\begin{proof}
The dual cost of one unit of $x_t(i)$ for the CES technology \eqref{eq:x_ces} is
\eqref{eq:px_dual}. Conditional on $x_t(i)$, Shephard's lemma gives
\eqref{eq:m_demand} and \eqref{eq:d_demand}. The Cobb-Douglas production function
\eqref{eq:production} has the unit cost \eqref{eq:unit_cost}; cost shares are $1-\alpha$ for labor
and $\alpha$ for the intermediate composite, which gives \eqref{eq:n_demand} and
\eqref{eq:x_demand}.
\end{proof}

\begin{assumption}[Bottleneck-relevant price region]\label{ass:input_saving}
The bottleneck-relevant region of the state space satisfies
\begin{equation}
    \frac{\tilde p_t^m}{p_t^d}
    \ge
    \frac{\omega(A_t)}{1-\omega(A_t)}.
    \label{eq:bottleneck_price_region}
\end{equation}
In this region, the effective imported input is sufficiently costly relative to the non-critical
substitute input.
\end{assumption}

\begin{lemma}[Adaptation is input-saving]\label{lem:input_saving}
Under Assumptions~\ref{ass:omega} and \ref{ass:input_saving}, adaptation weakly lowers the unit
cost of the intermediate bundle:
\begin{equation}
    \frac{\partial p_t^x(A,\tilde p_t^m,p_t^d)}{\partial A}\le 0.
    \label{eq:px_A_negative}
\end{equation}
It also weakly lowers imported-input use per unit of the intermediate composite:
\begin{equation}
    \frac{\partial}{\partial A}
    \left[
        \omega(A)^\rho
        \left(\frac{p_t^x(A,\tilde p_t^m,p_t^d)}{\tilde p_t^m}\right)^\rho
    \right]
    \le 0.
    \label{eq:m_per_x_A_negative}
\end{equation}
\end{lemma}

Thus, in the bottleneck-relevant region, an increase in installed adaptation lowers the firm's
imported-input intensity and weakly lowers marginal cost through the intermediate-input price. If
the imported-input cap binds, this also relaxes current scarcity pressure by lowering desired use
of the constrained input at a given production scale. Adaptation does not raise output directly; it
changes the cost-minimizing composition of inputs in favor of the substitute input.

\begin{proof}
Let
\[
    F(A)
    =
    \omega(A)^\rho(\tilde p_t^m)^{1-\rho}
    +
    (1-\omega(A))^\rho(p_t^d)^{1-\rho},
    \qquad
    p_t^x=F(A)^{1/(1-\rho)}.
\]
Differentiating gives
\[
    \frac{\partial \log p_t^x}{\partial A}
    =
    \frac{\rho\omega_A(A)}{1-\rho}
    \frac{
        \omega(A)^{\rho-1}(\tilde p_t^m)^{1-\rho}
        -
        (1-\omega(A))^{\rho-1}(p_t^d)^{1-\rho}
    }{F(A)}.
\]
For $\rho>1$ and for $\rho<1$, condition \eqref{eq:bottleneck_price_region} is equivalent to the
sign restriction needed for the right-hand side to be nonpositive when $\omega_A(A)<0$. Hence
\eqref{eq:px_A_negative} follows. For imported-input use per unit of $x$, use
\[
    \frac{m}{x}
    =
    \omega(A)^\rho
    \left(\frac{p_t^x(A,\tilde p_t^m,p_t^d)}{\tilde p_t^m}\right)^\rho .
\]
Its log derivative is
\[
    \rho\frac{\omega_A(A)}{\omega(A)}
    +
    \rho\frac{\partial \log p_t^x}{\partial A},
\]
which is weakly negative because $\omega_A(A)<0$ and
$\partial \log p_t^x/\partial A\le 0$.
\end{proof}

\begin{remark}[Input-saving region]
Assumption~\ref{ass:input_saving} restricts the region in which the baseline bottleneck mechanism
is studied. It does not introduce another disturbance. It says that adaptation is useful precisely
when the imported input is expensive or scarce relative to the available substitute. This rules out
a perverse case in which the firm spends on adaptation even though the imported input is cheap and
easy to obtain. The baseline is about repair under a bottleneck, so adaptation must save on the
input that has become costly or scarce.
\end{remark}

\subsection{Price Setting}

Firms set prices subject to Calvo frictions. In each period, a measure $1-\theta$ of firms can
reset their nominal price, while a measure $\theta\in(0,1)$ keeps its previous nominal price
unchanged. Let
\begin{equation}
    p_t^*=\frac{P_t^*}{P_t}
    \label{eq:reset_price_def}
\end{equation}
denote the reset price relative to the current aggregate price index. Aggregate gross inflation is
$\Pi_t=P_t/P_{t-1}$.

The price-index condition is
\begin{equation}
    1
    =
    (1-\theta)(p_t^*)^{1-\varepsilon}
    +
    \theta \Pi_t^{\varepsilon-1}.
    \label{eq:calvo_price_index}
\end{equation}
Price dispersion is
\begin{equation}
    \Delta_t
    =
    \int_0^1
    \left(\frac{P_t(i)}{P_t}\right)^{-\varepsilon}\dd i,
    \label{eq:price_dispersion_def}
\end{equation}
and evolves according to
\begin{equation}
    \Delta_t
    =
    (1-\theta)(p_t^*)^{-\varepsilon}
    +
    \theta \Pi_t^{\varepsilon}\Delta_{t-1}.
    \label{eq:price_dispersion_law}
\end{equation}

\begin{lemma}[Calvo reset-price condition]\label{lem:calvo_pricing}
Let real marginal cost be common across price-setting firms and equal to $\mc_t$. Define the
recursive pricing objects
\begin{align}
    \mathcal S_t^p
    &=
    \mc_tY_t
    +
    \theta\E_t M_{t,t+1}\Pi_{t+1}^{\varepsilon}\mathcal S_{t+1}^p,
    \label{eq:calvo_S}\\
    \mathcal F_t^p
    &=
    Y_t
    +
    \theta\E_t M_{t,t+1}\Pi_{t+1}^{\varepsilon-1}\mathcal F_{t+1}^p.
    \label{eq:calvo_F}
\end{align}
The optimal relative reset price satisfies
\begin{equation}
    p_t^*
    =
    \frac{\varepsilon}{\varepsilon-1}
    \frac{\mathcal S_t^p}{\mathcal F_t^p}.
    \label{eq:reset_price_foc}
\end{equation}
In a zero-inflation deterministic steady state with $\Pi_t=1$ and $\Delta_t=1$, the desired real
marginal cost is
\begin{equation}
    \bar{\mc}=\frac{\varepsilon-1}{\varepsilon}.
    \label{eq:ss_mc}
\end{equation}
\end{lemma}

The two pricing recursions have a simple interpretation. The numerator
$\mathcal S_t^p$ is the discounted marginal-cost burden expected over the life of a reset price,
while the denominator $\mathcal F_t^p$ is the corresponding discounted demand base. The reset
price is therefore a markup over expected marginal cost during the period in which the price may
remain fixed. A bottleneck matters for inflation because it raises the marginal-cost terms inside
$\mathcal S_t^p$, and monetary policy matters because it can change the expected demand and cost
path over the same pricing horizon.

\begin{proof}
A firm that resets in period $t$ chooses $P_t^*$ to maximize the expected discounted value of
profits while its price remains fixed. Conditional on not resetting for $k$ periods, its relative
price in period $t+k$ is
\[
    \frac{P_t^*}{P_{t+k}}
    =
    p_t^*
    \left(\Pi_{t+1}\Pi_{t+2}\cdots\Pi_{t+k}\right)^{-1}.
\]
Demand is isoelastic, so the first-order condition equates the reset price to a markup over the
discounted average of future marginal costs during the expected price duration. Collecting the
future marginal-cost terms into $\mathcal S_t^p$ and the future demand terms into
$\mathcal F_t^p$ gives \eqref{eq:calvo_S}--\eqref{eq:reset_price_foc}. In a deterministic
zero-inflation steady state, $p_t^*=1$ and the reset-price condition implies
$1=\varepsilon\bar{\mc}/(\varepsilon-1)$.
\end{proof}

Price dispersion is a real inefficiency. With common production technology and common input prices,
aggregate input requirements are proportional to $\Delta_tY_t$ rather than $Y_t$. Thus stabilizing
inflation matters not only because it stabilizes prices, but also because it limits inefficient
relative-price dispersion.

This is why the Calvo block is not cosmetic in the present model. When relative prices are
dispersed, the economy must produce more gross differentiated output to deliver the same final
output. Under an imported-input bottleneck, that extra gross production also increases pressure on
the scarce input. Price instability can therefore amplify the real resource cost of the bottleneck.

\section{Adaptation}

\subsection{Law of Motion and Cost}

The previous firm block describes how a given adaptation stock affects current costs. The dynamic
problem starts from the inverse question: how much should the firm spend today to change that stock
for tomorrow? Baseline adaptation evolves according to
\begin{equation}
    A_{t+1}=(1-\delta_A)A_t+I_t^A,
    \qquad
    \delta_A\in(0,1),
    \label{eq:A_law}
\end{equation}
where $I_t^A\ge 0$ is adaptation investment. The real installation cost is
\begin{equation}
    p_{A,t}\Psi(I_t^A),
    \qquad
    \Psi(I)=I+\frac{\phi_A}{2}I^2,
    \qquad
    \phi_A>0.
    \label{eq:adaptation_cost}
\end{equation}
The expenditure $p_{A,t}\Psi(I_t^A)$ represents the real resources used to install adaptation:
engineering time, testing, certification, retooling, supplier search, contract adjustment, and
process reorganization. These costs are incurred before the new production arrangement reduces
future imported-input needs, which is why adaptation is a forward-looking investment rather than a
static input choice.
The firm does not invest mechanically whenever the model contains the state variable $A_t$.
The state exists in every period, but positive investment occurs only when the bottleneck
environment makes additional input-saving capability valuable. The external states enter this
calculation through $p_t^m(D_t,X_t)$ and $\bar m_t(D_t,X_t)$: the first changes the external
procurement price, while the second can make the imported-input cap bind and generate the
scarcity rent $\chi_t$. Together they determine the effective price
$\tilde p_t^m=p_t^m+\chi_t$ that appears in current and expected future variable costs.
To allow monetary policy to affect the cost of front-loaded repair spending, suppose a fraction
$\vartheta_A\in[0,1]$ of the installation bill must be paid before current sales revenues are
received and is financed intraperiod at the gross policy rate $R_t$. The effective cost multiplier
is
\begin{equation}
    \Omega_t^A=(1-\vartheta_A)+\vartheta_A R_t.
    \label{eq:Omega_A}
\end{equation}
If $\vartheta_A=0$, adaptation is paid out of current resources and the direct financing-cost
channel is shut down. If $\vartheta_A>0$, tighter policy raises the current effective price of
installation.

This specification does not assume a full financial-frictions model. It is a working-capital-style
reduced form for front-loaded repair spending. If firms must pay engineers, certification costs,
equipment suppliers, or retooling costs before the new production arrangement generates revenues,
a higher policy rate raises the effective current cost of doing that repair. The multiplier
$\Omega_t^A$ should therefore be read as a real effective funding-cost wedge associated with the
gross policy rate, with no separate intraperiod inflation correction inside the installation bill.
This direct cost channel is combined below with the continuation value of repair, which depends on
future production scale and future scarcity.

\subsection{Firm Value and Adaptation FOC}

Calvo price dispersion creates price-history heterogeneity across product lines. To keep the
baseline focused on the bottleneck-repair mechanism, adaptation is chosen at the firm-technology
level and is common across the firm's price cohorts, in parallel with the standard Calvo assumption
that all price cohorts share the same marginal cost. The resource implication of price dispersion
is captured by replacing final output $Y_t$ with the gross production requirement $\Delta_tY_t$ in
variable costs.

\begin{assumption}[Common adaptation]\label{ass:common_adaptation}
Installed adaptation is a technology-level state shared by all price cohorts of a firm. A firm can
have different relative prices across cohorts because of Calvo frictions, but the same
$A_t$ governs the input requirement of every unit it produces.
\end{assumption}

Assumption~\ref{ass:common_adaptation} keeps price rigidity and repair dynamics separable. Calvo
frictions create dispersion in relative prices and therefore in output across price cohorts, but
they do not create separate adaptation stocks for different price histories. The effect of price
dispersion on the adaptation problem therefore enters through the gross production requirement
$\Delta_tY_t$. The reason is economic, not only technical: adaptation changes compatibility,
supplier structure, engineering standards, and the production process itself, so it is naturally
shared across price cohorts even when those cohorts charge different relative prices. A model in
which each price-age cohort carries its own adaptation stock would mix two heterogeneity problems,
price histories and repair histories; the baseline keeps one common repair technology so that the
bottleneck-repair mechanism remains the object being identified.

Because adaptation is common across the firm's Calvo price cohorts, the baseline writes the repair
problem directly at the symmetric-equilibrium technology level. This avoids introducing separate
adaptation states for each price history. In that representation, the real operating surplus of
the goods-producing sector before adaptation costs can be written as
\begin{equation}
    \pi_t^{op}
    =
    Y_t
    -
    \C_t^\Delta(Y_t,A_t;\tilde p_t^m,p_t^d,w_t,Z_t,\Delta_t),
    \label{eq:operating_profit}
\end{equation}
where the dispersion-adjusted variable-cost function is
\begin{equation}
    \C_t^\Delta(Y,A;\tilde p^m,p^d,w,Z,\Delta)
    =
    \mc_t(A,\tilde p^m,p^d,w,Z)\Delta Y.
    \label{eq:cost_function}
\end{equation}
The term $\Delta Y$ is the amount of gross differentiated output that must be produced to deliver
$Y$ units of final output when relative prices are dispersed.

This cost object is where the external bottleneck enters the firm's dynamic repair problem. Current
operating surplus depends on the current effective imported-input price
$\tilde p_t^m=p_t^m(D_t,X_t)+\chi_t$. Future continuation values depend on the future effective
price $\tilde p_{t+1}^m=p_{t+1}^m(D_{t+1},X_{t+1})+\chi_{t+1}$ and on future output and price
dispersion. Therefore the value of adaptation is high precisely when current or expected future
bottlenecks make imported-input dependence costly. In a tranquil environment with cheap,
unconstrained access and low expected future bottleneck risk, the same optimization can deliver the
corner solution $I_t^A=0$.

This object is not meant to be a separate sectoral planner problem. It is the normalized payoff for
the symmetric firm technology after imposing common adaptation and aggregating over Calvo price
cohorts. Individual firms still take aggregate prices, the scarcity rent, the Calvo pricing
objects, and the law of motion for aggregate states as given. They internalize how installed
adaptation lowers their future variable costs, but they do not internalize the effect of individual
imported-input demand on the aggregate scarcity rent. In symmetric equilibrium
$A_t(i)=A_t$, $I_t^A(i)=I_t^A$, and $\mc_t(i)=\mc_t$ for all firms with the same installed
technology. Operating surplus falls when the bottleneck raises $\tilde p_t^m$, when price
dispersion raises the gross production requirement, or when adaptation is too low to economize on
the imported input; the current-period gain from avoiding these losses is the force that makes
repair valuable.

The symmetric-equilibrium technology-level adaptation problem is therefore
\begin{equation}
    V_t(A_t;s_t)
    =
    \max_{\{I_t^A,A_{t+1}\}}
    \left\{
        \pi_t^{op}
        -
        \Omega_t^A p_{A,t}\Psi(I_t^A)
        +
        \E_t M_{t,t+1}
        V_{t+1}(A_{t+1};s_{t+1})
    \right\}
    \label{eq:firm_value}
\end{equation}
subject to \eqref{eq:A_law} and $I_t^A\ge 0$. The vector $s_t$ collects aggregate states and prices
taken as given by the individual firm.

\begin{definition}[Value of installed next-period adaptation]\label{def:QA}
Let
\begin{equation}
    \V_t^A
    =
    \frac{\partial V_t(A_t;s_t)}{\partial A_t}
    \label{eq:VA_def}
\end{equation}
be the marginal value of installed current adaptation capital. Define the date-$t$ value of one
unit of installed next-period adaptation as
\begin{equation}
    Q_t^A
    =
    \E_t M_{t,t+1}\V_{t+1}^A.
    \label{eq:QA_def}
\end{equation}
\end{definition}

\begin{proposition}[Adaptation Euler condition]\label{prop:adaptation_foc}
Optimal adaptation satisfies
\begin{equation}
    \Omega_t^A p_{A,t}\Psi'(I_t^A) - Q_t^A\ge 0,
    \qquad
    I_t^A\ge 0,
    \qquad
    I_t^A
    \left[
        \Omega_t^A p_{A,t}\Psi'(I_t^A) - Q_t^A
    \right]=0.
    \label{eq:adaptation_KKT}
\end{equation}
If the interior solution obtains, then
\begin{equation}
    I_t^A
    =
    \frac{1}{\phi_A}
    \left(
        \frac{Q_t^A}{\Omega_t^A p_{A,t}}-1
    \right).
    \label{eq:adaptation_policy_interior}
\end{equation}
With the nonnegativity constraint,
\begin{equation}
    I_t^A
    =
    \frac{1}{\phi_A}
    \left(
        \frac{Q_t^A}{\Omega_t^A p_{A,t}}-1
    \right)_+.
    \label{eq:adaptation_policy}
\end{equation}
The firm invests up to the point where the current marginal cost of one more unit of adaptation
equals the value of the future cost reductions that the unit creates. If the value $Q_t^A$ is below
the first unit's effective installation cost, the nonnegativity constraint binds and the firm does
not invest.
\end{proposition}

\begin{proof}
The choice of $I_t^A$ affects current payouts through
$-\Omega_t^A p_{A,t}\Psi(I_t^A)$ and affects the continuation value through
$A_{t+1}=(1-\delta_A)A_t+I_t^A$. The derivative of the continuation term with respect to
$I_t^A$ is $Q_t^A$ by Definition~\ref{def:QA}. The Kuhn-Tucker conditions for
$I_t^A\ge 0$ give \eqref{eq:adaptation_KKT}. Since $\Psi'(I)=1+\phi_A I$,
\eqref{eq:adaptation_policy_interior} and \eqref{eq:adaptation_policy} follow.
\end{proof}

\begin{proposition}[Continuation value of adaptation]\label{prop:QA_recursion}
The value of installed next-period adaptation satisfies
\begin{equation}
    Q_t^A
    =
    \E_t M_{t,t+1}
    \left[
        -\C_{A,t+1}^\Delta
        +
        (1-\delta_A)Q_{t+1}^A
    \right],
    \label{eq:QA_recursion}
\end{equation}
where
\begin{equation}
    \C_{A,t+1}^\Delta
    =
    \frac{\partial
    \C_{t+1}^\Delta(Y_{t+1},A_{t+1};
    \tilde p_{t+1}^m,p_{t+1}^d,w_{t+1},Z_{t+1},\Delta_{t+1})}
    {\partial A_{t+1}}.
    \label{eq:CA_def}
\end{equation}
Under Lemma~\ref{lem:input_saving}, $-\C_{A,t+1}^\Delta\ge 0$.
\end{proposition}

\begin{proof}
Differentiate \eqref{eq:firm_value} with respect to installed current adaptation $A_t$, using
the envelope theorem. Current adaptation affects current value through the variable-cost function
and affects next-period installed adaptation through
$A_{t+1}=(1-\delta_A)A_t+I_t^A$. Thus
\begin{equation}
    \V_t^A
    =
    -\C_{A,t}^\Delta+(1-\delta_A)Q_t^A.
    \label{eq:VA_envelope}
\end{equation}
Substitute \eqref{eq:VA_envelope} at date $t+1$ into Definition~\ref{def:QA}. Since
$\C_t^\Delta(Y,A;\cdot)=\mc_t(A,\cdot)\Delta_tY_t$ and Lemma~\ref{lem:input_saving} implies
$\partial\mc_t/\partial A\le 0$ in the bottleneck-relevant region, $-\C_{A,t+1}^\Delta\ge 0$.
\end{proof}

Equation \eqref{eq:QA_recursion} is the precise object behind the phrase ``value of repair.'' A
unit of adaptation is valuable because it lowers future production costs by reducing reliance on
the effective imported-input price $\tilde p_{t+1}^m=p_{t+1}^m+\chi_{t+1}$, and because installed
adaptation is persistent.

Define the one-period cost saving generated by one additional unit of installed adaptation as
\begin{equation}
    B_t^A
    \equiv
    -\C_{A,t}^\Delta
    \ge 0.
    \label{eq:adaptation_flow_benefit}
\end{equation}
The flow benefit $B_t^A$ is not an independent payoff term. It is the reduction in variable
production cost generated by a marginal increase in installed input-saving capability. It is large
when the firm is producing at scale, when price dispersion raises gross production requirements,
or when the effective imported-input price $\tilde p_t^m$ is high. It is small when production is
weak or when the imported input is cheap and unconstrained.
Under the transversality condition
\begin{equation}
    \lim_{K\to\infty}
    \E_t
    \left[
        M_{t,t+K}(1-\delta_A)^K Q_{t+K}^A
    \right]
    =
    0,
    \label{eq:adaptation_TVC}
\end{equation}
where $M_{t,t+k}\equiv\prod_{\ell=0}^{k-1}M_{t+\ell,t+\ell+1}$, iteration of
\eqref{eq:QA_recursion} gives
\begin{equation}
    Q_t^A
    =
    \sum_{k=1}^{\infty}
    (1-\delta_A)^{k-1}
    \E_t
    \left[
        M_{t,t+k}B_{t+k}^A
    \right].
    \label{eq:QA_forward}
\end{equation}
Thus $Q_t^A$ is the discounted value of future reductions in variable production costs generated
by one extra unit of installed adaptation: it is high when firms expect future bottlenecks to remain
tight, future imported-input rents to be high, or future production scale to be large, and it is low
when the firm expects the bottleneck to disappear quickly or future production to be weak. This
state-contingent value is the object through which monetary policy can affect recovery incentives.

\section{Symmetric Equilibrium}

\subsection{Aggregate Accounting}

With Calvo price dispersion, final output $Y_t$ requires $\Delta_tY_t$ units of gross
differentiated production. Let aggregate intermediate-composite use, imported-input use, and
aggregate non-critical substitute-input use be
\begin{align}
    \mathcal X_t
    &=
    \alpha\frac{\mc_t\Delta_tY_t}{p_t^x(A_t,\tilde p_t^m,p_t^d)},
    \label{eq:aggregate_x_calvo}\\
    M_t
    &=
    \mathcal X_t
    \omega(A_t)^\rho
    \left(
        \frac{p_t^x(A_t,\tilde p_t^m,p_t^d)}{\tilde p_t^m}
    \right)^\rho,
    \label{eq:aggregate_M_calvo}\\
    S_t
    &=
    \mathcal X_t
    \bigl(1-\omega(A_t)\bigr)^\rho
    \left(
        \frac{p_t^x(A_t,\tilde p_t^m,p_t^d)}{p_t^d}
    \right)^\rho.
    \label{eq:aggregate_S_calvo}
\end{align}
Aggregate labor demand is
\begin{equation}
    N_t
    =
    (1-\alpha)\frac{\mc_t\Delta_tY_t}{w_t}.
    \label{eq:aggregate_N_calvo}
\end{equation}
Aggregate real adaptation installation costs are
\begin{equation}
    \Theta_t^A
    =
    p_{A,t}\Psi(I_t^A).
    \label{eq:aggregate_adaptation_cost}
\end{equation}
The corresponding private adaptation payment is
\begin{equation}
    \Theta_t^{A,priv}
    =
    \Omega_t^A p_{A,t}\Psi(I_t^A).
    \label{eq:private_adaptation_cost}
\end{equation}
The aggregate accounting separates real absorption from private payments. The procurement cost
$p_t^mM_t$ is a real drain associated with obtaining the imported input from abroad. The
substitute-input cost $p_t^dS_t$ is the real cost of using non-critical alternatives inside the
baseline reduced-form supply block. The adaptation installation cost $\Theta_t^A$ is the real
resource cost of producing and installing repair services. By contrast, the scarcity payment
$\chi_tM_t$ is the value of permits or rents generated by the imported-input quantity ceiling and
is rebated to households through profits or transfers. The additional financing component
\[
    (\Omega_t^A-1)p_{A,t}\Psi(I_t^A)
\]
is a payment to lenders or households, not an additional absorption of real goods, unless one adds
a separate costly-intermediation technology.

In aggregate notation, the imported-input complementarity system is
\begin{equation}
    \chi_t\ge 0,
    \qquad
    \bar m_t-M_t\ge 0,
    \qquad
    \chi_t(\bar m_t-M_t)=0.
    \label{eq:aggregate_m_complementarity}
\end{equation}
Equation \eqref{eq:aggregate_m_complementarity} is the equilibrium version of
\eqref{eq:m_complementarity} after aggregating firms' cost-minimizing imported-input demands.

The aggregate resource constraint is therefore
\begin{equation}
    Y_t
    =
    C_t
    +
    p_t^m M_t
    +
    p_t^d S_t
    +
    \Theta_t^A.
    \label{eq:resource_constraint}
\end{equation}
Equation \eqref{eq:resource_constraint} keeps the distinction between the external procurement
cost $p_t^mM_t$, the domestic scarcity transfer $\chi_tM_t$, and the private financing payment
embedded in $\Theta_t^{A,priv}$. Firms take $p_t^m+\chi_t$ as the private marginal cost of
imported-input use and $\Omega_t^Ap_{A,t}$ as the private marginal installation cost of adaptation,
but only $p_t^mM_t$ and $p_{A,t}\Psi(I_t^A)$ absorb real resources at the aggregate level.

\begin{remark}[Private wedges and aggregate resource costs]\label{rem:private_wedges_transfer}
The terms $\chi_tM_t$ and $(\Omega_t^A-1)p_{A,t}\Psi(I_t^A)$ enter private firm payments but do not
enter the aggregate resource constraint. This is an accounting distinction rather than an
additional optimality condition.
The imported-input cap creates a shadow value for scarce units of the imported input. A firm that
uses one unit internalizes the opportunity cost $p_t^m+\chi_t$: $p_t^m$ is paid to procure the
input from abroad, while $\chi_t$ is the domestic rent attached to access to the constrained input.
At the aggregate level, the payment $p_t^mM_t$ is a real resource cost of obtaining imported goods.
The payment $\chi_tM_t$ is a transfer to the owner of the scarce entitlement or permit and is
rebated to households through profits or transfers. The same accounting logic applies to the
financing component of adaptation. The real installation resource cost is
$p_{A,t}\Psi(I_t^A)$; the additional payment
$(\Omega_t^A-1)p_{A,t}\Psi(I_t^A)$ is a transfer to the provider of working-capital finance. Hence
both wedges affect private marginal conditions but cancel in aggregate resource accounting.
\end{remark}

The distinction is economically useful. The scarcity rent affects inflation because firms price off
private marginal cost, and the financing wedge affects adaptation because firms invest off private
installation cost. Neither wedge is a technological destruction of resources in the baseline. The
external procurement cost $p_t^mM_t$, the substitute-input cost $p_t^dS_t$, and the real adaptation
installation cost $p_{A,t}\Psi(I_t^A)$ are the resource costs.

\begin{definition}[Symmetric recursive equilibrium]\label{def:equilibrium}
A symmetric recursive equilibrium collects the private decisions and market-clearing prices derived
above. It includes household choices $\{C_t,N_t,B_t^h\}$, firm choices
$\{Y_t,M_t,S_t,A_{t+1},I_t^A,\Pi_t,\Delta_t\}$, endogenous equilibrium prices
$\{w_t,\chi_t,R_t\}$, exogenous or state-implied input prices
$\{p_t^d,p_{A,t},p_t^m\}$, pricing objects $\{p_t^*,\mathcal S_t^p,\mathcal F_t^p\}$, and laws of
motion for the states. These objects form an equilibrium when:
\begin{enumerate}[leftmargin=1.5em]
    \item households satisfy \eqref{eq:hh_muc}--\eqref{eq:hh_euler};
    \item final demand satisfies \eqref{eq:variety_demand}--\eqref{eq:price_index};
    \item firms satisfy static cost minimization
    \eqref{eq:n_demand}--\eqref{eq:d_demand}, aggregate input requirements
    \eqref{eq:aggregate_x_calvo}--\eqref{eq:aggregate_N_calvo}, the Calvo pricing block
    \eqref{eq:calvo_price_index}--\eqref{eq:reset_price_foc}, and the adaptation conditions
    \eqref{eq:adaptation_KKT}--\eqref{eq:QA_recursion};
    \item the imported-input market satisfies
    \eqref{eq:m_cap}--\eqref{eq:m_complementarity}, equivalently
    \eqref{eq:aggregate_m_complementarity} in aggregate notation;
    \item the policy instrument $R_t$ is supplied by an admissible policy environment;
    \item aggregate accounting, labor-market clearing, asset-market clearing, and rebate accounting
    are satisfied:
    \begin{align}
        Y_t
        &=
        C_t+p_t^mM_t+p_t^dS_t+\Theta_t^A,
        \label{eq:resource_clearing}\\
        N_t
        &=
        (1-\alpha)\frac{\mc_t\Delta_tY_t}{w_t},
        \label{eq:labor_clearing}\\
        B_t^h&=0,
        \label{eq:bond_clearing}\\
        \frac{\Pi_t^f+T_t}{P_t}
        &=
        \pi_t^{op}
        -
        \Theta_t^{A,priv}
        +
        \chi_tM_t
        +
        (\Omega_t^A-1)p_{A,t}\Psi(I_t^A).
        \label{eq:rebate_clearing}
    \end{align}
    In \eqref{eq:labor_clearing}, the left-hand side is household labor supply and the right-hand
    side is aggregate firm labor demand, so the equation is labor-market clearing rather than a
    separate firm-only identity. Equation \eqref{eq:rebate_clearing} states one
    consolidated rebate accounting: operating profits are paid to the household, firms pay the
    private financed adaptation cost, and the scarcity-rent and financing components are returned
    to the representative household as in Assumption~\ref{ass:rent_rebate} and
    Remark~\ref{rem:private_wedges_transfer}. Since
    $\Theta_t^{A,priv}=p_{A,t}\Psi(I_t^A)+(\Omega_t^A-1)p_{A,t}\Psi(I_t^A)$, the adaptation terms
    in \eqref{eq:rebate_clearing} net out to the real installation cost
    $-p_{A,t}\Psi(I_t^A)$ at the household level. The financing wedge affects firms' private
    investment condition, but it is not an additional aggregate absorption of goods.
    Equivalently, household-side net income equals operating profits minus the real installation
    cost, with the scarcity-rent and financing wedges rebated back as transfers. Thus the household
    ultimately bears the real installation cost, while the scarcity and financing wedges net out
    through rebates and transfers.
\end{enumerate}
\end{definition}

For the model's mechanism, the central market-clearing condition is the imported-input
complementarity system \eqref{eq:m_complementarity}. A later resource-based repair extension can
replace the reduced-form substitute-input and adaptation-cost accounting with explicit domestic
labor and capital requirements. Thus the baseline deliberately stops before modeling the domestic
production side of substitutes and repair services. It allows firms to spend real resources on
substitution capability, but it does not yet ask whether those resources come from scarce domestic
engineers, specialized machinery, or workers reallocated away from current production. That
additional resource-allocation channel is kept out of the baseline so that the first mechanism is
clean: external access conditions create procurement costs and scarcity rents, while firm-level
adaptation changes future dependence on the constrained imported input.

\subsection{Equilibrium System}

Combining the household, firm, imported-input, pricing, adaptation, and policy blocks gives the
baseline equilibrium system. The household conditions determine demand, labor supply, and the
discount factor:
\begin{align}
    \Lambda_t&=C_t^{-\sigma},
    &
    N_t^\varphi&=w_t\Lambda_t,
    \label{eq:system_household_static}\\
    1&=\beta R_t\E_t\left[
        \frac{\Lambda_{t+1}}{\Lambda_t}\frac{1}{\Pi_{t+1}}
    \right].
    \label{eq:system_household_euler}
\end{align}
Given demand and discounting, the imported-input market determines whether the physical cap binds
and, if it binds, the scarcity rent that firms internalize:
\begin{align}
    \tilde p_t^m&=p_t^m+\chi_t,
    \label{eq:system_effective_price}\\
    \chi_t&\ge0,
    \qquad
    \bar m_t-M_t\ge0,
    \qquad
    \chi_t(\bar m_t-M_t)=0.
    \label{eq:system_input_market}
\end{align}
Firms' cost-minimizing input requirements then translate output, price dispersion, and effective
input prices into aggregate demands for labor, imported input, and substitute input:
\begin{align}
    \mathcal X_t
    &=
    \alpha\frac{\mc_t\Delta_tY_t}{p_t^x(A_t,\tilde p_t^m,p_t^d)},
    \label{eq:system_x}\\
    M_t
    &=
    \mathcal X_t\omega(A_t)^\rho
    \left(\frac{p_t^x(A_t,\tilde p_t^m,p_t^d)}{\tilde p_t^m}\right)^\rho,
    \label{eq:system_M}\\
    S_t
    &=
    \mathcal X_t(1-\omega(A_t))^\rho
    \left(\frac{p_t^x(A_t,\tilde p_t^m,p_t^d)}{p_t^d}\right)^\rho,
    \label{eq:system_S}\\
    N_t
    &=
    (1-\alpha)\frac{\mc_t\Delta_tY_t}{w_t}.
    \label{eq:system_labor}
\end{align}
The Calvo block maps marginal costs into reset prices, inflation, and the law of motion for price
dispersion:
\begin{align}
    1
    &=
    (1-\theta)(p_t^*)^{1-\varepsilon}
    +
    \theta\Pi_t^{\varepsilon-1},
    \label{eq:system_calvo_price_index}\\
    \Delta_t
    &=
    (1-\theta)(p_t^*)^{-\varepsilon}
    +
    \theta\Pi_t^\varepsilon\Delta_{t-1},
    \label{eq:system_calvo_dispersion}\\
    \mathcal S_t^p
    &=
    \mc_tY_t
    +
    \theta\E_tM_{t,t+1}\Pi_{t+1}^{\varepsilon}\mathcal S_{t+1}^p,
    \label{eq:system_calvo_S}\\
    \mathcal F_t^p
    &=
    Y_t
    +
    \theta\E_tM_{t,t+1}\Pi_{t+1}^{\varepsilon-1}\mathcal F_{t+1}^p,
    \label{eq:system_calvo_F}\\
    p_t^*
    &=
    \frac{\varepsilon}{\varepsilon-1}
    \frac{\mathcal S_t^p}{\mathcal F_t^p}.
    \label{eq:system_reset_price}
\end{align}
The adaptation block links today's installation decision to the future value of lower
imported-input dependence:
\begin{align}
    A_{t+1}
    &=
    (1-\delta_A)A_t+I_t^A,
    \label{eq:system_A_law}\\
    \Omega_t^Ap_{A,t}\Psi'(I_t^A)-Q_t^A&\ge0,
    \qquad
    I_t^A\ge0,
    \label{eq:system_adaptation_kkt_1}\\
    I_t^A[\Omega_t^Ap_{A,t}\Psi'(I_t^A)-Q_t^A]&=0,
    \label{eq:system_adaptation_kkt}\\
    Q_t^A
    &=
    \E_tM_{t,t+1}
    [-\C_{A,t+1}^{\Delta}+(1-\delta_A)Q_{t+1}^A].
    \label{eq:system_Q}
\end{align}
Given the policy instrument $R_t$, aggregate resources close the private equilibrium system:
\begin{align}
    Y_t
    &=
    C_t+p_t^mM_t+p_t^dS_t
    +
    p_{A,t}\Psi(I_t^A).
    \label{eq:system_resource}
\end{align}
The exogenous states
\[
    (D_t,X_t,\lambda_t^D,\lambda_t^X,Z_t,p_t^d,p_{A,t})
\]
evolve according to the laws specified above. The system is organized in the same order as the
mechanism: the external block determines $p_t^m$ and $\bar m_t$, the market-clearing block
determines the endogenous scarcity rent $\chi_t$, and the dynamic firm block determines whether
firms spend resources today to reduce future dependence on the constrained imported input.

\begin{definition}[Recursive state vector]\label{def:recursive_state}
The baseline recursive state is
\begin{equation}
    \mathcal Z_t
    =
    (D_t,X_t,\lambda_t^D,\lambda_t^X,Z_t,p_t^d,p_{A,t},A_t,\Delta_{t-1}),
    \label{eq:recursive_state}
\end{equation}
with $\lambda_t^D$ and $\lambda_t^X$ omitted when event intensities are constant. The predetermined
state variables are the external bottleneck stocks, the event intensities, productivity, competitive
supply prices, installed adaptation, and inherited price dispersion. The equilibrium objects
\[
    (C_t,Y_t,N_t,M_t,S_t,\chi_t,\Pi_t,p_t^*,\Delta_t,I_t^A,A_{t+1},
    \mathcal S_t^p,\mathcal F_t^p,Q_t^A)
\]
are time-$t$ policy functions of $\mathcal Z_t$.
\end{definition}

Definition~\ref{def:recursive_state} makes the private baseline economy directly usable in
recursive policy problems. The external states determine the current imported-input price and cap;
$A_t$ determines current input intensity; and $\Delta_{t-1}$ carries the real effects of past price
dispersion into current production requirements. A policy environment specified below completes
the mapping from states to the policy instrument $R_t$.

\begin{assumption}[Regular recursive equilibrium]\label{ass:regular_equilibrium}
The equilibrium considered below is a regular interior recursive equilibrium for household choices,
Calvo reset prices, and the relevant pricing objects; the imported-input constraint is allowed to
bind or be slack through the complementarity system. Value functions are finite, the household bond
transversality condition holds, and the adaptation transversality condition
\eqref{eq:adaptation_TVC} holds.

More precisely, for any fixed active set of complementarity constraints in the local neighborhood
considered below, the smooth part of the equilibrium system has full local rank with respect to the
contemporaneous equilibrium objects. Hence, once the policy environment specifies how $R_t$ is
determined, the baseline system locally defines recursive policy functions for the private
allocation and prices. This assumption is not a global existence theorem. It is the local regularity
condition needed for the comparative statics and policy problems in the next part.
\end{assumption}

\part{Monetary Policy in the Baseline Economy}

\section{Policy Closure and Flexible-Price Benchmark}

The baseline economy above defines the private equilibrium constraints: households, firms,
imported-input market clearing, Calvo price setting, and adaptation. The monetary-policy problem is
introduced only after those constraints are in place. This separation is deliberate. The first part
defines how bottlenecks and repair work; the second part asks how a policy instrument interacts
with current inflation, current scarcity, and future repair incentives.

\begin{definition}[Interest-rate notation]\label{def:interest_rate_notation}
The monetary-policy instrument is the gross nominal one-period rate
\begin{equation}
    R_t=1+i_t,
    \label{eq:R_nominal_definition}
\end{equation}
where $i_t$ is the net nominal policy rate. The corresponding ex post gross real return between
$t$ and $t+1$ is $R_t/\Pi_{t+1}$. This real return is not denoted by $R_t$.

Given a flexible-price benchmark allocation, let
\begin{equation}
    \Lambda_t^n=(C_t^n)^{-\sigma}
    \label{eq:natural_muc}
\end{equation}
denote benchmark marginal utility. The gross natural real rate is
\begin{equation}
    \mathcal R_t^n
    =
    \left[
        \beta
        \E_t
        \left(
            \frac{\Lambda_{t+1}^n}{\Lambda_t^n}
        \right)
    \right]^{-1}.
    \label{eq:gross_natural_real_rate}
\end{equation}
The associated gross natural nominal intercept at inflation target $\bar\Pi$ is
\begin{equation}
    R_t^n=\bar\Pi\,\mathcal R_t^n,
    \qquad
    r_t^n=\mathcal R_t^n-1.
    \label{eq:natural_nominal_intercept}
\end{equation}
Thus $R_t$ is the actual gross nominal policy rate, $\mathcal R_t^n$ is the gross natural real
rate, and $r_t^n$ is the net natural real rate. The notation keeps separate the policy instrument,
the realized real return, and the flexible-price natural-rate benchmark.
\end{definition}

The same convention is used throughout the policy analysis: symbols without a superscript denote
equilibrium objects under the current policy environment, while the superscript $n$ denotes the
flexible-price benchmark used only to define natural output and the natural-rate intercept. The
aggregate imported input $M_t$ is the realized quantity after the availability cap is imposed,
whereas $M_t^d(\cdot)$ below denotes firms' desired demand before imposing the cap. The real
adaptation resource cost is $\Theta_t^A=p_{A,t}\Psi(I_t^A)$, while
$\Theta_t^{A,priv}=\Omega_t^Ap_{A,t}\Psi(I_t^A)$ is the private financed payment entering the
firm's investment condition. These distinctions keep the policy instrument, natural benchmark,
scarcity rent, and repair cost from being collapsed into a single reduced-form wedge.

The working-capital term $\Omega_t^A=(1-\vartheta_A)+\vartheta_A R_t$ uses the actual gross
nominal policy rate because it is a reduced-form funding-cost channel for current adaptation
expenditure. It is not a natural-rate object. A bottleneck-adjusted Taylor rule can later replace
the fixed intercept $\bar R$ with the state-contingent nominal benchmark $R_t^n$, but the local
mechanism below first uses a standard Taylor-type closure to identify how a policy tightening
propagates through the bottleneck economy.

In this local experiment the central bank does not change the external bottleneck itself. The
stocks $D_t$ and $X_t$ are held fixed, so policy does not reopen trade routes, remove sanctions, or
increase the physical input cap directly. Instead, the policy rate works through private choices:
it changes current absorption through the household Euler equation, the current scarcity rent
through imported-input demand, and repair incentives through the private cost and continuation
value of adaptation.

\begin{definition}[Flexible-price benchmark]\label{def:flex_benchmark}
Fix the current production state
\[
    s_t^n=(D_t,X_t,Z_t,p_t^d,A_t).
\]
The flexible-price benchmark is the static production allocation obtained by setting
\begin{equation}
    \theta=0,
    \qquad
    p_t^*=1,
    \qquad
    \Delta_t=1,
    \qquad
    \mc_t^n=\frac{\varepsilon-1}{\varepsilon},
    \label{eq:flex_price_conditions}
\end{equation}
while keeping the same household preferences, production technology, imported-input constraint,
external procurement price $p_t^m$, available imported-input capacity $\bar m_t$, productivity
$Z_t$, substitute-input price $p_t^d$, and installed adaptation $A_t$. The benchmark output
$Y_t^n$ is the final-good output that satisfies the static cost-minimization conditions, the
imported-input complementarity system, household intratemporal optimality, and the static goods
resource constraint under \eqref{eq:flex_price_conditions}.
\end{definition}

Constructively, the benchmark is pinned down by the same static household labor condition, the
same cost-minimizing input demands, the same imported-input market-clearing complementarity
system, and the static goods resource constraint
\begin{equation}
    Y_t^n
    =
    C_t^n
    +
    p_t^mM_t^n
    +
    p_t^dS_t^n,
    \label{eq:flex_static_resource}
\end{equation}
evaluated under flexible prices and zero price dispersion. Current adaptation expenditure is not
included in \eqref{eq:flex_static_resource}; it is a dynamic repair choice compared separately in
the policy analysis, not part of the static natural-output benchmark.

\begin{assumption}[Local well-defined flexible-price benchmark]\label{ass:flex_benchmark_unique}
For every state in the local neighborhood considered below, the system defining
$Y_t^n$ in Definition~\ref{def:flex_benchmark} has a locally unique solution.
\end{assumption}

The flexible-price benchmark is not a separate economy and does not remove the imported-input
bottleneck. It removes only nominal price rigidity and relative-price dispersion. Therefore
$Y_t/Y_t^n$ measures the output gap relative to the allocation attainable under the same real
external bottleneck and the same installed adaptation stock. A negative output gap is therefore not
the full real damage from the external disruption: the benchmark already embeds the real bottleneck,
and the gap measures the additional inefficiency generated by nominal rigidity and policy
conditions relative to that constrained real benchmark.

This is why the benchmark is useful in a bottleneck economy. It holds fixed the real scarcity
environment and asks what output would be attainable if firms could adjust prices without Calvo
frictions. The output gap therefore separates real scarcity from the additional distortions created
by sticky prices and policy-induced demand conditions.

The same benchmark also preserves the intertemporal household channel familiar from persistent
supply-regime models. Because the flexible-price allocation depends on the bottleneck state, the
benchmark path of consumption and marginal utility changes with expected disruption and relief.
The natural real rate $\mathcal R_t^n$ therefore remains state dependent. In the present model,
however, this natural-rate channel is not the only object of interest. It provides the benchmark
for policy rules and the discounting of future repair payoffs, while the distinctive mechanism is
the interaction between current scarcity and firms' endogenous adaptation.

The benchmark is current-state and production-based. It conditions on installed adaptation $A_t$,
because $A_t$ is predetermined at the start of period $t$, and it does not use the current
adaptation choice $I_t^A$ to redefine natural output. Current repair spending affects aggregate
resources and future exposure through the firm's dynamic problem, but including it in the benchmark
would make the output-gap object mechanically depend on the repair decision whose policy effects are
being studied. The output-gap term in the policy rule therefore isolates the contemporaneous nominal
wedge relative to the same real bottleneck and installed technology.

The baseline policy rule is a Taylor-type interest-rate rule:
\begin{equation}
    R_t
    =
    \bar R
    \left(\frac{\Pi_t}{\bar\Pi}\right)^{\phi_\Pi}
    \left(\frac{Y_t}{Y_t^n}\right)^{\phi_Y}
    \exp(\varepsilon_t^R),
    \qquad
    \phi_\Pi>1,\quad \phi_Y\ge 0,
    \label{eq:taylor_rule}
\end{equation}
where $Y_t^n$ is the flexible-price benchmark associated with the same current production state and
installed adaptation stock, and $\varepsilon_t^R$ is an exogenous monetary-policy innovation. The comparative
statics below consider a contractionary innovation, i.e. a local increase in $R_t$ induced by
$\varepsilon_t^R$ holding the external bottleneck state fixed.

\begin{remark}[Policy innovation versus rule change]\label{rem:policy_innovation}
The derivative with respect to $R_t$ used below is a local derivative with respect to the
monetary-policy innovation $\varepsilon_t^R$ in \eqref{eq:taylor_rule}. It is not a permanent
change in the coefficients of the policy rule. Thus the comparative statics hold the external
states $(D_t,X_t)$, installed adaptation $A_t$, and the rule parameters
$(\bar R,\phi_\Pi,\phi_Y)$ fixed. They describe a tighter current policy stance in the same
external bottleneck environment, not a different long-run policy regime or a different stochastic
process for bottlenecks.
\end{remark}

\section{Policy Environments}

The rule \eqref{eq:taylor_rule} is useful for identifying the local transmission mechanism, but it
is not the only policy environment of interest. The baseline economy can be closed in four ways:
a fixed-intercept Taylor rule, a bottleneck-adjusted Taylor rule, optimal policy under discretion,
and optimal policy under commitment. This mirrors the logic of regime-dependent policy analysis in
persistent-supply-regime models, but the state dependence here is continuous and comes from the
bottleneck economy rather than from a discrete normal/bad regime.

The sequence is intentional. The two rules ask whether the policy intercept moves with the
state-dependent natural benchmark. Discretion and commitment ask a deeper question: once the
bottleneck economy makes repair endogenous, what allocation would a policymaker choose subject to
the private equilibrium constraints? A rule can correct the natural-rate intercept and still leave
the future repair margin outside the object being optimized.

The four environments therefore differ along three dimensions. The fixed Taylor rule keeps a
constant intercept and reacts only to realized inflation and the output gap. The
bottleneck-adjusted rule keeps the same feedback structure but replaces the constant intercept with
the natural benchmark implied by the current bottleneck state. Discretion drops the mechanical rule
and lets the central bank choose an implementable allocation period by period, taking future
re-optimization as given. Commitment goes one step further: the central bank chooses an
intertemporal implementable plan, so promises about future policy become part of the current
equilibrium through expectations and the value of repair.

\subsection{Implementable Equilibrium Constraints}

For policy analysis, the long private equilibrium system is compressed into a small number of
implementability objects. This is only notation: it does not add a new agent or a new constraint.
Collect the period-$t$ private equilibrium objects in
\begin{equation}
    \mathbf y_t
    =
    (C_t,Y_t,N_t,M_t,S_t,\chi_t,\Pi_t,p_t^*,\Delta_t,I_t^A,A_{t+1},
    \mathcal S_t^p,\mathcal F_t^p,Q_t^A,R_t).
    \label{eq:policy_choice_vector}
\end{equation}
The recursive state is $\mathcal Z_t$ from Definition~\ref{def:recursive_state}. Let
\begin{equation}
    \mathcal G^B
    \left(
        \mathbf y_t,
        \E_t\mathbf y_{t+1};
        \mathcal Z_t
    \right)
    =
    0
    \label{eq:baseline_implementability}
\end{equation}
denote the system formed by the household conditions, static firm conditions, Calvo pricing
recursions, adaptation Euler equation, imported-input market clearing, aggregate accounting, and
the laws of motion for predetermined states. Let
\begin{equation}
    \mathcal K^B
    \left(
        \mathbf y_t,
        \E_t\mathbf y_{t+1};
        \mathcal Z_t
    \right)
    \ge 0,
    \qquad
    \boldsymbol\eta_t^B\ge 0,
    \qquad
    \boldsymbol\eta_t^B\circ \mathcal K^B_t=0
    \label{eq:baseline_complementarity_set}
\end{equation}
collect the complementarity restrictions, including the imported-input cap and the non-negativity
condition on adaptation investment. The superscript $B$ indicates the baseline economy. Later
extensions will replace $\mathcal G^B$ and $\mathcal K^B$ with the corresponding extended
private-equilibrium constraints.

This notation does not hide a different economic problem. It is only a compact representation of
the equations derived above. The system $\mathcal G^B=0$ collects the smooth restrictions that make
an allocation privately implementable: demand must satisfy the household Euler equation, prices
must satisfy Calvo optimality, and repair must satisfy the firm's adaptation Euler condition. The
system $\mathcal K^B\ge0$ collects the non-smooth restrictions: the imported-input cap may bind or
not bind, and adaptation investment may be positive or at a corner. The central bank therefore does
not choose arbitrary allocations; it chooses objects that can be decentralized by households and
firms in the bottleneck economy. The policy rate $R_t$ is included because it is required to satisfy
the household Euler condition and, when $\vartheta_A>0$, the adaptation-cost implementability
condition. It is not an independent primitive state or an ad hoc control outside the equilibrium
system; it is the instrument through which a candidate allocation is implemented.

The role of the policy rate in the implementability system can be seen directly from two equations.
First, the household Euler equation contains the actual gross nominal rate,
\begin{equation}
    \mathcal G_t^{E}
    \equiv
    1-\beta R_t
    \E_t
    \left[
        \frac{\Lambda_{t+1}}{\Lambda_t}
        \frac{1}{\Pi_{t+1}}
    \right]
    =
    0.
    \label{eq:policy_euler_constraint}
\end{equation}
Second, on the interior adaptation branch, the firm's repair condition contains the same policy
rate through the effective funding cost,
\begin{equation}
    \mathcal G_t^{A}
    \equiv
    \Omega_t^Ap_{A,t}\Psi'(I_t^A)-Q_t^A
    =
    0,
    \qquad
    \Omega_t^A=(1-\vartheta_A)+\vartheta_AR_t.
    \label{eq:policy_adaptation_constraint}
\end{equation}
Thus $R_t$ is not merely a label attached to a policy rule. It appears in the Euler equation that
implements intertemporal demand and in the repair condition that implements the firm's adaptation
decision. Locally,
\begin{equation}
    \frac{\partial \mathcal G_t^{E}}{\partial R_t}
    =
    -\beta
    \E_t
    \left[
        \frac{\Lambda_{t+1}}{\Lambda_t}
        \frac{1}{\Pi_{t+1}}
    \right],
    \qquad
    \frac{\partial \mathcal G_t^{A}}{\partial R_t}
    =
    \vartheta_Ap_{A,t}\Psi'(I_t^A)\ge0.
    \label{eq:R_direct_implementability_derivatives}
\end{equation}
The first derivative is the demand-implementation margin; the second is the direct repair-cost
margin. Current adaptation also has a resource and a state component:
\begin{equation}
    \frac{\partial}{\partial I_t^A}
    \left[p_{A,t}\Psi(I_t^A)\right]
    =
    p_{A,t}\Psi'(I_t^A),
    \qquad
    \frac{\partial A_{t+1}}{\partial I_t^A}=1.
    \label{eq:I_resource_state_margins}
\end{equation}
A policymaker therefore cannot evaluate repair only as a current expenditure. The same unit of
$I_t^A$ absorbs resources today and changes the next state that determines future dependence on the
constrained imported input.

Inflation enters the policy problem through the Calvo implementability conditions rather than as an
ad hoc loss term. The gross inflation rate $\Pi_t$ appears in the price-index condition and in the
law of motion for price dispersion,
\begin{equation}
    \Delta_t
    =
    (1-\theta)(p_t^*)^{-\varepsilon}
    +
    \theta\Pi_t^\varepsilon\Delta_{t-1}.
    \label{eq:policy_dispersion_margin}
\end{equation}
Because aggregate input requirements are proportional to $\Delta_tY_t$, inflation affects welfare
through the implementability of sticky prices and through the real resource cost of price
dispersion. This is why the optimal-policy problems below choose implementable allocations rather
than imposing a reduced-form quadratic loss in inflation.

\subsection{Rule-Based Policy}

The fixed-intercept Taylor rule is
\begin{equation}
    R_t^{F}
    =
    \bar R
    \left(\frac{\Pi_t}{\bar\Pi}\right)^{\phi_\Pi}
    \left(\frac{Y_t}{Y_t^n}\right)^{\phi_Y}
    \exp(\varepsilon_t^R).
    \label{eq:fixed_taylor_rule}
\end{equation}
This rule is state-blind in the same sense as a fixed-intercept Taylor rule in a regime-switching
model: the intercept does not move with the current bottleneck state. It reacts to realized
inflation and the output gap, but it does not directly adjust to how $(D_t,X_t,A_t)$ change the
flexible-price benchmark and the natural real rate. Hence it can confuse two objects that the
model keeps separate: a nominal deviation from target and a real deterioration in the constrained
flexible-price allocation. When the bottleneck worsens, the natural benchmark itself can move; a
fixed intercept treats that state-dependent real adjustment as if the same nominal stance remained
appropriate.

The bottleneck-adjusted Taylor rule replaces the fixed intercept with the natural nominal
intercept from Definition~\ref{def:interest_rate_notation}:
\begin{equation}
    R_t^{BA}
    =
    R_t^n
    \left(\frac{\Pi_t}{\bar\Pi}\right)^{\phi_\Pi}
    \left(\frac{Y_t}{Y_t^n}\right)^{\phi_Y}
    \exp(\varepsilon_t^R).
    \label{eq:bottleneck_adjusted_taylor_rule}
\end{equation}
Here $R_t^n=\bar\Pi\mathcal R_t^n$ is not a discretionary policy choice; it is the nominal
counterpart of the flexible-price natural real rate implied by the same bottleneck environment.
Thus the rule adjusts to the state-dependent natural benchmark generated by the current external
disruption, relief, and installed adaptation stock.


The rule-based comparison separates two questions. The first is the Galo--Nu\~no-type question:
whether a fixed intercept creates a nominal misalignment when the natural benchmark is
state-dependent. The second is specific to this model: whether correcting that natural-rate
misalignment is enough when current policy also affects future repair incentives. The
bottleneck-adjusted rule can remove one source of policy error, namely the failure to track the
state-dependent natural rate. It does not, by construction, choose $R_t$ by comparing the value of
lower current inflation with the value of faster future substitution away from the constrained
input.

\subsection{Optimal Policy Under Discretion}

Under discretion, the central bank re-optimizes each period and takes the future policy functions
used by private agents as given. The objective is household welfare,
\begin{equation}
    U(C_t,N_t)
    =
    \frac{C_t^{1-\sigma}}{1-\sigma}
    -
    \frac{N_t^{1+\varphi}}{1+\varphi}.
    \label{eq:period_welfare}
\end{equation}
Let $V^D(\mathcal Z)$ denote the value of the discretionary policy problem. A Markov discretionary
equilibrium solves
\begin{equation}
    V^D(\mathcal Z_t)
    =
    \max_{\mathbf y_t}
    \left\{
        U(C_t,N_t)
        +
        \beta
        \E_t V^D(\mathcal Z_{t+1})
    \right\}
    \label{eq:discretion_bellman}
\end{equation}
subject to
\begin{equation}
    \mathcal G^B
    \left(
        \mathbf y_t,
        \E_t\mathbf Y^D(\mathcal Z_{t+1});
        \mathcal Z_t
    \right)
    =
    0,
    \qquad
    \mathcal K^B
    \left(
        \mathbf y_t,
        \E_t\mathbf Y^D(\mathcal Z_{t+1});
        \mathcal Z_t
    \right)
    \ge 0,
    \label{eq:discretion_constraints}
\end{equation}
where $\mathbf Y^D(\cdot)$ is the future discretionary policy function anticipated by households
and firms. The choice variable $\mathbf y_t$ includes the implementation rate $R_t$, but the
central bank is not free to pick $R_t$ independently of private behavior. The rate must support an
allocation satisfying the household Euler equation, the pricing block, and the adaptation Euler
condition.

This makes discretion a current-state policy problem. The central bank internalizes that a higher
$R_t$ can lower current demand, reduce pressure on the imported-input cap, and affect inflation
through marginal costs. It also internalizes that today's policy changes next period's
predetermined states, especially installed adaptation $A_{t+1}$ and price dispersion $\Delta_t$.
What it cannot do is bind future choices: households and firms evaluate future pricing and repair
using the future discretionary policy function $\mathbf Y^D(\cdot)$, not a promise made today.

For interior constraints, the discretionary Lagrangian is
\begin{equation}
    \mathcal L_t^D
    =
    U(C_t,N_t)
    +
    \beta\E_t V^D(\mathcal Z_{t+1})
    +
    (\boldsymbol\mu_t^D)'
    \mathcal G^B
    \left(
        \mathbf y_t,
        \E_t\mathbf Y^D(\mathcal Z_{t+1});
        \mathcal Z_t
    \right),
    \label{eq:discretion_lagrangian}
\end{equation}
with complementary slackness conditions for $\mathcal K^B$. A discretionary solution satisfies,
for every interior component $a_t$ of $\mathbf y_t$,
\begin{equation}
    \frac{\partial U(C_t,N_t)}{\partial a_t}
    +
    \beta
    \E_t
    \left[
        V^D_{\mathcal Z}(\mathcal Z_{t+1})
        \frac{\partial \mathcal Z_{t+1}}{\partial a_t}
    \right]
    +
    (\boldsymbol\mu_t^D)'
    \frac{\partial \mathcal G_t^B}{\partial a_t}
    =
    0,
    \label{eq:discretion_foc_generic}
\end{equation}
plus the private implementability constraints. This is the formal sense in which the central bank
solves a problem under discretion: it chooses an implementable equilibrium allocation taking
future discretionary behavior as given. It internalizes the current welfare effects and the effect
of today's choices on predetermined states such as $A_{t+1}$ and $\Delta_t$, but it does not bind
future policy through promises.

The generic first-order condition \eqref{eq:discretion_foc_generic} therefore has a direct
economic reading. The derivative of $\mathcal G_t^B$ prices the current implementability
constraints: the chosen allocation must remain consistent with Euler equations, Calvo pricing,
input-market clearing, and repair optimality. The derivative of $\mathcal Z_{t+1}$ prices the
fact that current policy changes tomorrow's state. This is where repair enters the discretionary
policy problem: a policy-induced decline in current adaptation changes $A_{t+1}$ and therefore
future exposure to the bottleneck.

For the implementation rate itself, the direct discretionary margin is especially transparent. Let
$\mu_t^{D,E}$ denote the multiplier on the household Euler implementability condition
\eqref{eq:policy_euler_constraint} and let $\mu_t^{D,A}$ denote the multiplier on the interior
adaptation condition \eqref{eq:policy_adaptation_constraint}. Abstracting from constraints in which
$R_t$ enters only indirectly through other equilibrium variables, the direct part of the
discretionary first-order condition for $R_t$ is
\begin{equation}
    \mu_t^{D,E}
    \left[
        -\beta
        \E_t
        \left(
            \frac{\Lambda_{t+1}}{\Lambda_t}
            \frac{1}{\Pi_{t+1}}
        \right)
    \right]
    +
    \mu_t^{D,A}
    \vartheta_Ap_{A,t}\Psi'(I_t^A)
    =
    0.
    \label{eq:discretion_R_direct_margin}
\end{equation}
This expression is not a new behavioral equation for households or firms; it is the policymaker's
shadow-value comparison. The first term is the value of using the nominal rate to implement
intertemporal demand. The second term is the value of the same rate change through the private cost
of repair. In equilibrium, the central bank cannot move one margin without moving the other because
both are supported by the same instrument $R_t$. Appendix~\ref{app:computational_residual_system}
then states the corresponding DEQN residual system used for training; in the discretionary case the
value-gradient terms are obtained from the differentiable value network rather than from a separate
value-gradient network.

\subsection{Optimal Policy Under Commitment}

Under commitment, the central bank chooses an implementable plan and can influence private
expectations through promises about future policy. This is a stronger policy environment because
the value of repair is forward-looking. If firms expect future policy to keep output weak or
future bottlenecks less valuable to escape, the continuation value $Q_t^A$ falls and current
adaptation becomes less attractive. If promises about future policy make repair more valuable,
current investment can respond before those future states arrive. The Ramsey problem is
\begin{equation}
    \max_{\{\mathbf y_t\}_{t\ge 0}}
    \E_0
    \sum_{t=0}^{\infty}
    \beta^t
    U(C_t,N_t)
    \label{eq:commitment_ramsey_problem}
\end{equation}
subject to, for all $t$,
\begin{equation}
    \mathcal G^B
    \left(
        \mathbf y_t,
        \E_t\mathbf y_{t+1};
        \mathcal Z_t
    \right)
    =
    0,
    \qquad
    \mathcal K^B
    \left(
        \mathbf y_t,
        \E_t\mathbf y_{t+1};
        \mathcal Z_t
    \right)
    \ge 0,
    \label{eq:commitment_constraints}
\end{equation}
and the exogenous and endogenous state laws. The associated Lagrangian is
\begin{equation}
    \mathcal L^C
    =
    \E_0
    \sum_{t=0}^{\infty}
    \beta^t
    \left[
        U(C_t,N_t)
        +
        (\boldsymbol\mu_t^C)'
        \mathcal G^B
        \left(
            \mathbf y_t,
            \E_t\mathbf y_{t+1};
            \mathcal Z_t
        \right)
    \right],
    \label{eq:commitment_lagrangian}
\end{equation}
again with complementary slackness for $\mathcal K^B$.

The commitment first-order condition for an interior component $a_t$ of $\mathbf y_t$ has the
generic current-value form
\begin{equation}
    \frac{\partial U(C_t,N_t)}{\partial a_t}
    +
    (\boldsymbol\mu_t^C)'
    \frac{\partial \mathcal G_t^B}{\partial a_t}
    +
    \boldsymbol\varrho_{t-1}^C(a_t)
    +
    \E_t
    \left[
        \beta
        (\boldsymbol\mu_{t+1}^C)'
        \frac{\partial \mathcal G_{t+1}^B}{\partial a_t}
    \right]
    =
    0.
    \label{eq:commitment_foc_generic}
\end{equation}
The term $\boldsymbol\varrho_{t-1}^C(a_t)$ collects the normalized lagged promise terms attached to
date-$t$ objects that entered date-$t-1$ forward-looking private constraints. The last term captures
the effect of $a_t$ on future implementability through predetermined states. These two terms are the
formal differences from discretion. In this model, the forward-looking private constraints include
the household Euler equation, the Calvo pricing recursions, and the adaptation-value recursion for
$Q_t^A$. Commitment therefore works not only by shaping inflation expectations, but also by shaping
the expected value of repair.

This is the reason commitment is substantively important in the baseline economy. In a standard
sticky-price model, commitment matters because current promises affect inflation expectations and
the output gap. Here the same logic also applies to structural repair. A promise about future
policy changes the expected path of demand, scarcity rents, and marginal costs; those objects
enter $Q_t^A$ and therefore affect the firm's current willingness to invest in lower future
imported-input dependence.

The direct current-period margin for $R_t$ contains the same two current terms as
\eqref{eq:discretion_R_direct_margin}, with Ramsey multipliers replacing discretionary multipliers.
The difference is not that commitment invents a different current instrument. The difference is that
FOCs for forward-looking or state-shifting objects also contain the lagged promise term and the
future-state term in \eqref{eq:commitment_foc_generic}. Thus commitment can affect current inflation
and repair through the expected future path of implementability constraints, not only through the
contemporaneous Euler and adaptation-cost margins.

In a recursive implementation, commitment generally augments the physical state $\mathcal Z_t$ with
lagged promise variables attached to forward-looking private constraints. The promise variables are
defined in scaled form. Let
\[
    \mathbf p_{t+1}^C
    =
    \mathcal P(\boldsymbol\mu_t^C,\mathbf y_t)
\]
denote the vector of inherited promises selected from the date-$t$ Ramsey multipliers and scaled
by the date-$t$ marginal-utility normalizers that appear in the forward-looking constraints. For
example, a raw lagged pricing multiplier that would appear in the next-period FOC multiplied by
$C_t^\sigma C_{t+1}^{-\sigma}$ is carried forward as a scaled object proportional to
$\mu_t C_t^\sigma$. This is only a normalization. Current consumption $C_t$ remains a control in
period $t$; the state carries the inherited scaled promise, not current consumption itself. With
this normalization, the compact recursive representation is
\begin{equation}
    \mathcal Z_t^C
    =
    \left(
        \mathcal Z_t,
        \mathbf p_t^{E},
        \mathbf p_t^{P},
        \mathbf p_t^{A}
    \right),
    \label{eq:commitment_state}
\end{equation}
where $\mathbf p_t^{E}$ summarizes scaled promises attached to the household Euler equation,
$\mathbf p_t^{P}$ summarizes scaled promises attached to the Calvo pricing block, and
$\mathbf p_t^{A}$ summarizes scaled promises attached to the adaptation-value recursion. This is
the same economic state augmentation as carrying raw lagged multipliers together with the lagged
normalizing variables, but it avoids putting lagged consumption into the state vector as a separate
object. The exact dimension of these promise states is determined by the implementation of
\eqref{eq:commitment_constraints}. The expression in \eqref{eq:commitment_state} should therefore
be read as the recursive architecture of the commitment problem, not as a claim that all numerical
promise-state dimensions have already been fixed. The economic role is clear: commitment makes
future policy promises state variables because private agents price and invest based on expected
future policy.

The recursive state used by each policy environment follows from this distinction:
\begin{equation}
    \mathcal Z_t^{F}
    =
    \mathcal Z_t^{BA}
    =
    \mathcal Z_t^{D}
    =
    \mathcal Z_t,
    \qquad
    \mathcal Z_t^{C}
    =
    \left(
        \mathcal Z_t,
        \mathbf p_t^{E},
        \mathbf p_t^{P},
        \mathbf p_t^{A}
    \right).
    \label{eq:policy_environment_states}
\end{equation}
The two rule-based environments and discretion are Markov in the physical bottleneck state. A
recursive implementation of commitment is Markov in the physical state plus the relevant promise
variables. This is the formal difference between re-optimizing every period and choosing a
history-dependent plan.

\subsection{Policy Comparison in the Baseline Economy}

The four policy environments answer different but nested questions. The comparison between the
fixed Taylor rule and the bottleneck-adjusted Taylor rule isolates the role of state dependence in
the policy intercept: it asks whether policy performs poorly simply because it fails to track the
natural benchmark generated by the bottleneck state. The comparison between the adjusted rule and
discretion isolates the value of optimization: it asks whether a formula that tracks the natural
rate is enough, or whether the central bank must explicitly choose the current inflation--scarcity--
repair trade-off. The comparison between discretion and commitment isolates the value of promises:
it asks whether managing expectations about future policy changes current price setting and current
adaptation investment.

The comparison is deliberately stronger than a rule-only exercise. A bottleneck-adjusted Taylor
rule may track the natural real rate and stabilize the nominal side better than a fixed rule, but
it need not optimize the intertemporal repair margin created by the policy-sensitive adaptation
decision. Optimal policy is the benchmark for that intertemporal margin. The central
question for the quantitative part of the dissertation is therefore whether natural-rate adjustment
is sufficient once supply recovery is endogenous, or whether discretion and commitment prescribe a
different policy path because today's stance changes tomorrow's repair capacity.

For each policy environment $P\in\{F,BA,D,C\}$, the objects to be compared are the policy function
for the rate, the induced allocation, and the bottleneck-repair outcomes:
\begin{equation}
    R_t^P(\mathcal Z_t^P),
    \qquad
    \mathcal O_t^P
    =
    \left(
        \Pi_t^P,
        \frac{Y_t^P}{Y_t^n(\mathcal Z_t^P)},
        \chi_t^P,
        I_t^{A,P},
        A_{t+1}^P
    \right).
    \label{eq:policy_comparison_outcomes}
\end{equation}
The first two components are the conventional nominal and real-stabilization objects. The third
component measures the current scarcity price of access to the constrained imported input. The last
two components measure structural repair: current adaptation expenditure and the next-period stock
of installed input-saving capability. These objects keep the policy comparison aligned with the
model's mechanism rather than reducing it to inflation alone. The notation
$Y_t^n(\mathcal Z_t^P)$ means the same flexible-price benchmark evaluated at the physical state
induced by policy $P$; the benchmark concept itself is not policy-specific.
The vector also separates a current symptom from a future state. The rent $\chi_t^P$ measures how
tight the imported-input allocation problem is today, while $I_t^{A,P}$ and $A_{t+1}^P$ measure
whether firms are becoming less dependent on that allocation problem tomorrow. A policy can look
successful on the current rent and still leave the economy with weaker future repair.

Welfare comparisons use the household objective implied by the same equilibrium allocation:
\begin{equation}
    W^P(\mathcal Z_0^P)
    =
    \E_0
    \sum_{t=0}^{\infty}
    \beta^t
    U(C_t^P,N_t^P).
    \label{eq:policy_welfare}
\end{equation}
When a consumption-equivalent comparison is useful, the gain of policy $P$ relative to a reference
policy $Q$ is the scalar $\lambda^{P,Q}$ satisfying
\begin{equation}
    \E_0
    \sum_{t=0}^{\infty}
    \beta^t
    U\!\left((1+\lambda^{P,Q})C_t^Q,N_t^Q\right)
    =
    W^P(\mathcal Z_0^P).
    \label{eq:consumption_equivalent}
\end{equation}
This welfare object is not a substitute for the mechanism-level comparison in
\eqref{eq:policy_comparison_outcomes}. It summarizes the household value of the induced allocation,
while the outcome vector identifies whether the welfare difference comes with lower inflation,
lower current scarcity, faster repair, or some combination of these margins.

The mechanism-level policy experiment combines two maps. The external bottleneck map is
\[
    (D_t,X_t)
    \longrightarrow
    (p_t^m,\bar m_t)
    \longrightarrow
    \chi_t
    \longrightarrow
    (\mc_t,\Pi_t,Y_t),
\]
where disruption and relief determine procurement costs and physical availability, and a binding
availability constraint generates the endogenous scarcity rent. The monetary-policy map starts
from the same external state and asks how a tighter current policy stance changes the equilibrium
objects generated by that bottleneck:
\[
    R_t\uparrow
    \quad\Longrightarrow\quad
    \underbrace{\Pi_t\downarrow,\ \chi_t\downarrow}_{\text{current stabilization and bottleneck relief}}
    \qquad\text{but}\qquad
    \underbrace{I_t^A\downarrow,\ A_{t+1}\downarrow}_{\text{slower structural repair}}.
\]
The first map explains why the shock matters; the second map explains where the monetary trade-off
comes from. The local results below state the conditions under which each arrow has the displayed
sign.

\section{Local Conditional Equilibrium Implications}

The results in this section are local sufficient sign results around a regular recursive
equilibrium. They are not global monotonicity theorems and should not be read as claiming that a
monetary tightening lowers inflation or repair incentives in every possible equilibrium. The
economic reason is straightforward. Tightening lowers current demand and therefore tends to reduce
imported-input demand and current scarcity pressure. Inflation falls only if this effect lowers
expected Calvo marginal costs overall. Adaptation slows only if the decline in expected repair
benefits, together with any direct funding-cost effect, dominates offsetting forces. The signs for
inflation and for the value of repair therefore require explicit equilibrium restrictions on
expected marginal costs and expected future repair benefits. The signs are derived in the same
order as the economic mechanism. First, external
disruption and relief determine whether the effective imported-input cost is driven by procurement
prices, scarcity rents, or both. Second, conditional on a binding bottleneck, monetary tightening
reduces current pressure on the constrained input by lowering output demand. Third, the Calvo block
maps the same marginal-cost pressure into inflation. Finally, the adaptation Euler equation shows
how the policy stance affects the investment that reduces future input dependence.

\begin{definition}[Desired imported-input demand]\label{def:desired_import_demand}
For a given current state and effective imported-input price, let $M_t^d(\cdot)$ denote the
unconstrained aggregate imported-input demand implied by firms' static cost-minimization conditions
before imposing the aggregate availability cap. It is the amount firms would like to use at the
given prices and production scale. Actual equilibrium use $M_t$ is obtained only after the
complementarity system with $\bar m_t$ and $\chi_t$ is imposed. In local derivatives, $M_t^d(\cdot)$
is the demand generated by \eqref{eq:aggregate_M_calvo} at the relevant equilibrium values of the
remaining arguments. When only the effective price varies, write $M_t^d(\tilde p_t^m)$ and define
$M_p=\partial M_t^d/\partial\tilde p_t^m$. When output and the effective price both vary, write
$M_t^d(Y_t,\tilde p_t^m,A_t)$ and define $M_Y=\partial M_t^d/\partial Y_t$ and
$M_p=\partial M_t^d/\partial\tilde p_t^m$ evaluated at the local equilibrium.
\end{definition}

\subsection{External Bottleneck Tightness}

\begin{proposition}[Direct disruption, relief, and effective imported-input cost]\label{prop:DX_effects}
Suppose Assumption~\ref{ass:DX_mapping} holds. External disruption raises the procurement price
and lowers available imported-input supply:
\begin{equation}
    \frac{\partial p_t^m}{\partial D_t}=\nu_{pD}p_t^m\ge 0,
    \qquad
    \frac{\partial \bar m_t}{\partial D_t}=-\nu_{qD}\bar m_t\le 0.
    \label{eq:D_derivatives}
\end{equation}
External relief lowers the procurement price and raises available supply:
\begin{equation}
    \frac{\partial p_t^m}{\partial X_t}=-\nu_{pX}p_t^m\le 0,
    \qquad
    \frac{\partial \bar m_t}{\partial X_t}=\nu_{qX}\bar m_t\ge 0.
    \label{eq:X_derivatives}
\end{equation}
Holding current output, installed adaptation, price dispersion, productivity, wages, and the
substitute-input price fixed, if the imported-input constraint is locally binding and aggregate
desired imported-input demand is locally downward-sloping in $\tilde p_t^m$, then the direct
scarcity-rent responses are
\begin{align}
    \frac{\partial\chi_t}{\partial D_t}
    &=
    \frac{\nu_{qD}\bar m_t}{-M_p}
    -
    \nu_{pD}p_t^m,
    \label{eq:chi_D_exact}\\
    \frac{\partial\chi_t}{\partial X_t}
    &=
    -
    \frac{\nu_{qX}\bar m_t}{-M_p}
    +
    \nu_{pX}p_t^m,
    \label{eq:chi_X_exact}
\end{align}
where $M_p=\partial M_t^d/\partial\tilde p_t^m<0$. Hence disruption raises the scarcity rent if
the availability effect dominates the direct price-induced demand reduction:
\begin{equation}
    \frac{\nu_{qD}\bar m_t}{-M_p}
    \ge
    \nu_{pD}p_t^m,
    \label{eq:D_quantity_dominance}
\end{equation}
and relief lowers the scarcity rent if the availability improvement dominates the direct
price-induced demand increase:
\begin{equation}
    \frac{\nu_{qX}\bar m_t}{-M_p}
    \ge
    \nu_{pX}p_t^m.
    \label{eq:X_quantity_dominance}
\end{equation}
\end{proposition}

\begin{proof}
The derivatives in \eqref{eq:D_derivatives} and \eqref{eq:X_derivatives} follow directly from
\eqref{eq:pm_DX}--\eqref{eq:mbar_DX}. If the constraint binds, direct aggregate imported-input demand
$M_t^d(\tilde p_t^m)$ satisfies
\begin{equation}
    M_t^d(p_t^m+\chi_t)=\bar m_t.
    \label{eq:binding_m_market}
\end{equation}
Totally differentiating with respect to any external state $s\in\{D_t,X_t\}$ gives
\begin{equation}
    (M_t^d)'(\tilde p_t^m)
    \left(
        \frac{\partial p_t^m}{\partial s}
        +
        \frac{\partial\chi_t}{\partial s}
    \right)
    =
    \frac{\partial \bar m_t}{\partial s}.
    \label{eq:chi_DX_derivative_proof}
\end{equation}
Since $(M_t^d)'(\tilde p_t^m)<0$, the signs in
\eqref{eq:D_derivatives} and \eqref{eq:X_derivatives} imply
\eqref{eq:chi_D_exact}--\eqref{eq:chi_X_exact}. Conditions
\eqref{eq:D_quantity_dominance}--\eqref{eq:X_quantity_dominance} give the stated rent signs.
\end{proof}

The proposition separates the two direct components of the effective imported-input price:
\begin{equation}
    \tilde p_t^m
    =
    \underbrace{p_t^m}_{\text{external procurement cost}}
    +
    \underbrace{\chi_t}_{\text{endogenous scarcity rent}}.
    \label{eq:effective_price_decomp}
\end{equation}
This decomposition disciplines the mechanism. The model can generate inflation pressure from a
higher external import price, but the bottleneck mechanism is present only when the quantity
constraint produces a positive $\chi_t$. Equations
\eqref{eq:chi_D_exact}--\eqref{eq:chi_X_exact} explain why the price and quantity channels must be
kept conceptually separate in the direct effect. The terms
$\nu_{qD}\bar m_t/(-M_p)$ and $\nu_{qX}\bar m_t/(-M_p)$ measure how much the scarcity rent must move
to ration a smaller or larger physical input stock. The terms $\nu_{pD}p_t^m$ and
$\nu_{pX}p_t^m$ measure how much of that rationing is already done by the external procurement
price. A higher procurement price lowers desired imported-input demand and can mechanically reduce
the scarcity rent, while a tighter physical cap raises it. Thus worse external access is not
one-dimensional. If disruption mainly raises procurement costs, firms may reduce desired
imported-input use and the scarcity rent need not rise; if disruption mainly cuts physical
availability, the rent rises because the economy is trying to allocate a smaller input stock across
firms. The latter case is the core bottleneck case, and the rent is a genuine equilibrium object
rather than a relabeling of the external import price. General-equilibrium responses of output,
wages, and policy can add further effects; the proposition isolates the primitive external
price-versus-quantity mapping.

\subsection{Monetary Tightening and Current Scarcity}

The sign of the current scarcity response to monetary policy is a general-equilibrium object. It
requires the standard New Keynesian demand response: a contractionary policy innovation lowers
current absorption and output, all else equal.

\begin{lemma}[Input-demand slopes]\label{lem:local_M_slopes}
At fixed cost shifters and real marginal cost, aggregate cost-minimizing imported-input demand in
\eqref{eq:aggregate_M_calvo} is locally increasing in final output and locally decreasing in the
effective imported-input price:
\begin{equation}
    M_Y\equiv \frac{\partial M_t}{\partial Y_t}>0,
    \qquad
    M_p\equiv \frac{\partial M_t}{\partial \tilde p_t^m}<0.
    \label{eq:local_M_slopes}
\end{equation}
\end{lemma}

\begin{proof}
From \eqref{eq:aggregate_M_calvo},
\begin{equation}
    M_t
    =
    \alpha
    \frac{\mc_t\Delta_tY_t}{p_t^x(A_t,\tilde p_t^m,p_t^d)}
    \omega(A_t)^\rho
    \left(
        \frac{p_t^x(A_t,\tilde p_t^m,p_t^d)}{\tilde p_t^m}
    \right)^\rho .
    \label{eq:M_for_slope_proof}
\end{equation}
At fixed $(A_t,\Delta_t,w_t,p_t^d,Z_t,\mc_t)$, the expression is proportional to $Y_t$, so
$M_Y>0$. For the own-price derivative, rewrite \eqref{eq:M_for_slope_proof} as
\[
    M_t
    =
    K_t
    \left[p_t^x(A_t,\tilde p_t^m,p_t^d)\right]^{\rho-1}
    (\tilde p_t^m)^{-\rho},
    \qquad
    K_t=\alpha\mc_t\Delta_tY_t\omega(A_t)^\rho>0.
\]
Let
\[
    s_t^m
    \equiv
    \frac{\partial\log p_t^x}{\partial\log \tilde p_t^m}
    \in(0,1)
\]
denote the cost share of the imported input in the CES intermediate bundle. Then
\[
    \frac{\partial\log M_t}{\partial\log \tilde p_t^m}
    =
    (\rho-1)s_t^m-\rho
    =
    -\left[\rho-(\rho-1)s_t^m\right]
    <
    0,
\]
because $s_t^m\in(0,1)$ and $\rho>0$. Hence $M_p<0$.
\end{proof}

\begin{assumption}[Local demand response]\label{ass:demand_response}
At the equilibrium under consideration, a contractionary monetary-policy innovation satisfies
\begin{equation}
    \frac{\partial Y_t}{\partial R_t}<0
    \label{eq:Y_R_negative}
\end{equation}
holding the external states $(D_t,X_t)$ and installed adaptation $A_t$ fixed.
\end{assumption}

Assumption~\ref{ass:demand_response} is the standard New Keynesian demand-side restriction in this
local experiment. It is not a statement that the central bank directly changes the external
bottleneck. The external procurement price and the availability cap are held fixed; the policy rate
works by changing current absorption and production. This is why the scarcity response below is a
scale effect: lower output reduces desired use of the constrained imported input at the existing
technology and external input conditions.

\begin{proposition}[Current bottleneck relief]\label{prop:current_bottleneck_relief}
Suppose the imported-input constraint binds, Lemma~\ref{lem:local_M_slopes} applies locally, and
Assumption~\ref{ass:demand_response} holds. Then a contractionary monetary-policy innovation
weakly lowers the current scarcity rent:
\begin{equation}
    \frac{\partial \chi_t}{\partial R_t}\le 0.
    \label{eq:chi_R_negative}
\end{equation}
\end{proposition}

\begin{proof}
When the cap binds,
\begin{equation}
    M_t^d(Y_t,\tilde p_t^m,A_t)=\bar m_t.
    \label{eq:M_market_R}
\end{equation}
Holding $(D_t,X_t)$ and $A_t$ fixed, $\bar m_t$ and $p_t^m$ are fixed. Totally differentiating
with respect to $R_t$ gives
\begin{equation}
    M_Y\frac{\partial Y_t}{\partial R_t}
    +
    M_p\frac{\partial \chi_t}{\partial R_t}
    =
    0,
    \label{eq:chi_R_proof}
\end{equation}
where $M_Y>0$ and $M_p<0$ by Lemma~\ref{lem:local_M_slopes}. Since
$\partial Y_t/\partial R_t<0$, \eqref{eq:chi_R_proof} implies
$\partial\chi_t/\partial R_t\le 0$.
\end{proof}

Hence the current bottleneck-relief channel operates through quantities rather than repair: tight
policy reduces output demand, lower output requires less of the constrained imported input, and the
shadow price of the constraint falls. The production structure has not become less dependent on the
imported input; the economy is using less of it because it produces less today.

\subsection{Monetary Tightening and Calvo Inflation}

The inflation response is disciplined by the Calvo pricing block, not by a reduced-form inflation
equation. Around a zero-inflation deterministic reference allocation, equations
\eqref{eq:calvo_price_index}--\eqref{eq:reset_price_foc} imply the standard local New Keynesian
Phillips curve
\begin{equation}
    \widehat\pi_t
    =
    \beta\E_t\widehat\pi_{t+1}
    +
    \kappa_\Pi\widehat{\mc}_t,
    \qquad
    \kappa_\Pi
    =
    \frac{(1-\theta)(1-\beta\theta)}{\theta}
    >
    0,
    \label{eq:calvo_nkpc}
\end{equation}
where hats denote log deviations from the reference allocation and
$\widehat\pi_t=\log\Pi_t$. Price dispersion $\Delta_t$ remains a nonlinear state in the full
baseline model; \eqref{eq:calvo_nkpc} is used only to state the local sign of the
nominal-stabilization channel.

\begin{assumption}[Local Calvo marginal-cost response]\label{ass:calvo_mc_response}
At the equilibrium under consideration, a contractionary monetary-policy innovation weakly lowers
expected real marginal costs along the relevant Calvo pricing horizon:
\begin{equation}
    \frac{\partial}{\partial R_t}
    \E_t\widehat{\mc}_{t+j}
    \le 0,
    \qquad
    j=0,1,2,\ldots.
    \label{eq:mc_path_R_negative}
\end{equation}
\end{assumption}

\begin{proposition}[Conditional nominal stabilization under Calvo pricing]\label{prop:calvo_inflation_response}
Suppose the Calvo block is locally represented by \eqref{eq:calvo_nkpc}, the no-bubble forward
solution for inflation is selected, and Assumption~\ref{ass:calvo_mc_response} holds. Then a
contractionary monetary-policy innovation weakly lowers current inflation:
\begin{equation}
    \frac{\partial \Pi_t}{\partial R_t}\le 0.
    \label{eq:Pi_R_negative_calvo}
\end{equation}
\end{proposition}

\begin{proof}
Forward iteration of \eqref{eq:calvo_nkpc} gives
\begin{equation}
    \widehat\pi_t
    =
    \kappa_\Pi
    \sum_{j=0}^{\infty}
    \beta^j
    \E_t\widehat{\mc}_{t+j}.
    \label{eq:calvo_forward_solution}
\end{equation}
Since $\kappa_\Pi>0$ and Assumption~\ref{ass:calvo_mc_response} states that the contractionary
innovation weakly lowers each expected marginal-cost term, it weakly lowers
$\widehat\pi_t$. Because $\Pi_t=\exp(\widehat\pi_t)$ locally around the reference allocation, the
same sign applies to gross inflation.
\end{proof}

In this model, the nominal-stabilization force works through the same marginal-cost object that
contains ordinary production costs and the effective cost of the constrained imported input. The
proposition does not assume that the sign of the marginal-cost response is automatic. A tighter
policy stance reduces current demand, which tends to lower imported-input demand and scarcity
pressure, but other equilibrium prices may also move. The proposition therefore states the
conditional result: under the local Calvo marginal-cost response in
Assumption~\ref{ass:calvo_mc_response}, inflation falls.

\subsection{Monetary Tightening and Adaptation}

The current effect of monetary policy on adaptation can operate through two components. The first
is the direct installation-cost channel through $\Omega_t^A$. The second is the continuation-value
channel through $Q_t^A$: if a contractionary innovation lowers expected future production scale,
expected future price-dispersion costs, or expected future scarcity rents, it lowers the private
benefit of reducing imported-input dependence.

\begin{assumption}[Local repair-benefit response]\label{ass:Q_response}
At the equilibrium under consideration, a contractionary policy innovation weakly lowers the
expected discounted future cost savings generated by installed adaptation:
\begin{equation}
    \frac{\partial}{\partial R_t}
    \E_t
    \left[
        M_{t,t+k}B_{t+k}^A
    \right]
    \le 0,
    \qquad
    k=1,2,\ldots.
    \label{eq:future_repair_benefit_R_negative}
\end{equation}
\end{assumption}

The condition links the policy innovation to the firm's dynamic payoff from repair. The private
return to adaptation is high when future production is large, future price dispersion is costly, or
future imported-input scarcity is high. Tightening may lower this return by reducing expected
future scale and scarcity rents; it may also create offsetting effects by changing future costs.
Assumption~\ref{ass:Q_response} restricts attention to local equilibria in which the contractionary
policy innovation reduces discounted repair benefits overall. Without this restriction, the sign of
the continuation-value response is not a primitive implication of the baseline.

\begin{proposition}[Conditional continuation-value response]\label{prop:Q_response}
Under \eqref{eq:QA_forward} and Assumption~\ref{ass:Q_response}, a contractionary
monetary-policy innovation weakly lowers the continuation value of adaptation:
\begin{equation}
    \frac{\partial Q_t^A}{\partial R_t}\le 0.
    \label{eq:Q_R_negative}
\end{equation}
\end{proposition}

\begin{proof}
Differentiate \eqref{eq:QA_forward} term by term. Each term has nonnegative weight
$(1-\delta_A)^{k-1}$ and is weakly decreasing in $R_t$ by Assumption~\ref{ass:Q_response}. Hence
the sum is weakly decreasing.
\end{proof}

\begin{proposition}[Conditional adaptation slowdown]\label{prop:adaptation_slowdown}
Suppose the interior adaptation solution obtains, $\vartheta_A\ge 0$, and the conditions of
Proposition~\ref{prop:Q_response} hold. Then a contractionary monetary-policy innovation weakly
lowers current adaptation investment:
\begin{equation}
    \frac{\partial I_t^A}{\partial R_t}\le 0.
    \label{eq:I_R_negative}
\end{equation}
If either $\vartheta_A>0$ or $\partial Q_t^A/\partial R_t<0$, the inequality is strict.
\end{proposition}

\begin{proof}
From \eqref{eq:adaptation_policy_interior},
\begin{equation}
    I_t^A
    =
    \frac{1}{\phi_A}
    \left(
        \frac{Q_t^A}{\Omega_t^A p_{A,t}}-1
    \right).
    \label{eq:I_policy_for_derivative}
\end{equation}
At fixed $p_{A,t}$,
\begin{equation}
    \frac{\partial I_t^A}{\partial R_t}
    =
    \frac{1}{\phi_A}
    \left[
        \frac{1}{\Omega_t^A p_{A,t}}
        \frac{\partial Q_t^A}{\partial R_t}
        -
        \frac{Q_t^A}{p_{A,t}(\Omega_t^A)^2}
        \frac{\partial\Omega_t^A}{\partial R_t}
    \right].
    \label{eq:I_R_derivative}
\end{equation}
Since $\partial\Omega_t^A/\partial R_t=\vartheta_A\ge 0$, $Q_t^A\ge 0$ by
\eqref{eq:QA_forward}, and $\partial Q_t^A/\partial R_t\le 0$ by
Proposition~\ref{prop:Q_response}, the derivative is weakly negative.
\end{proof}

The structural-recovery channel follows from this Euler equation under the stated local
conditions. A higher policy rate can make current repair spending more expensive and, when
Assumption~\ref{ass:Q_response} holds, can lower the expected future benefit of repair. The firm
then invests less in reducing its future dependence on the constrained input.

\subsection{Baseline Trade-Off}

\begin{proposition}[Conditional baseline monetary-policy trade-off]\label{prop:baseline_tradeoff}
Suppose Assumptions~\ref{ass:demand_response}, \ref{ass:calvo_mc_response}, and
\ref{ass:Q_response} hold, the imported-input constraint binds, aggregate imported-input demand is
locally increasing in output and decreasing in the effective imported-input price, the Calvo
no-bubble inflation solution is selected, and the interior adaptation solution obtains. Then a
contractionary monetary-policy innovation has the following local effects:
\begin{align}
    \frac{\partial \Pi_t}{\partial R_t}&\le 0,
    \label{eq:tradeoff_pi}\\
    \frac{\partial \chi_t}{\partial R_t}&\le 0,
    \label{eq:tradeoff_chi}\\
    \frac{\partial I_t^A}{\partial R_t}&\le 0.
    \label{eq:tradeoff_I}
\end{align}
The second inequality is the direct bottleneck-relief result. The first and third inequalities are
conditional on the additional local restrictions in Assumptions~\ref{ass:calvo_mc_response} and
\ref{ass:Q_response}, respectively.
\end{proposition}

\begin{proof}
Proposition~\ref{prop:calvo_inflation_response} gives \eqref{eq:tradeoff_pi} from the Calvo
pricing block. Proposition~\ref{prop:current_bottleneck_relief} gives \eqref{eq:tradeoff_chi} from
the imported-input complementarity system. Proposition~\ref{prop:adaptation_slowdown} gives
\eqref{eq:tradeoff_I} from the adaptation Euler condition.
\end{proof}

Together, the three inequalities are a conditional local sign result, not an unconditional policy
ranking. Under the stated local restrictions, a contractionary innovation delivers one direct
current bottleneck-relief effect and two additional conditional responses: lower inflation through
the Calvo marginal-cost channel and lower repair through the adaptation Euler channel:
\begin{equation}
    R_t\uparrow
    \quad\Rightarrow\quad
    \Pi_t\downarrow,\qquad
    \chi_t\downarrow,\qquad
    I_t^A\downarrow.
    \label{eq:three_channels_summary}
\end{equation}
Under these sufficient local conditions, monetary tightening can stabilize current inflation and
relax the current bottleneck by lowering today's demand for the scarce input, while weakening the
investment that would make firms less exposed to that bottleneck tomorrow. The baseline trade-off
is therefore intertemporal: current stabilization versus structural escape.

\subsection{Limiting Cases}

The same equilibrium system nests three useful limiting cases. These cases are not alternative
models; they clarify which assumptions are doing work in the baseline mechanism.

\begin{corollary}[Slack imported-input cap]\label{cor:slack_cap}
If $\bar m_t>M_t$ locally, then $\chi_t=0$. The imported-input distortion reduces to the external
procurement price $p_t^m$, and the current bottleneck-relief channel through
$\partial\chi_t/\partial R_t$ is absent.
\end{corollary}

\begin{proof}
The result follows directly from the complementarity condition
\eqref{eq:aggregate_m_complementarity}. If the cap is slack, the multiplier on the cap is zero.
\end{proof}

\begin{corollary}[No active repair]\label{cor:no_repair}
If
\begin{equation}
    Q_t^A\le \Omega_t^Ap_{A,t},
    \label{eq:no_repair_condition}
\end{equation}
then $I_t^A=0$. In this region, monetary policy can still affect inflation and current scarcity,
but there is no active structural-recovery response in period $t$.
\end{corollary}

\begin{proof}
At $I_t^A=0$, $\Psi'(0)=1$. Condition \eqref{eq:no_repair_condition} is exactly the
Kuhn-Tucker condition \eqref{eq:adaptation_KKT} for the nonnegativity constraint to bind.
\end{proof}

\begin{corollary}[No direct financing-cost channel]\label{cor:no_direct_financing}
If $\vartheta_A=0$, then $\Omega_t^A=1$ and the direct effect of $R_t$ on the installation cost of
adaptation is shut down. Monetary policy can still affect repair through the continuation value
$Q_t^A$.
\end{corollary}

\begin{proof}
Set $\vartheta_A=0$ in \eqref{eq:Omega_A}. The adaptation policy
\eqref{eq:adaptation_policy} then has no direct $R_t$ term in the installation cost. Under
Assumption~\ref{ass:competitive_supply_prices}, the local policy derivative at fixed exogenous
$p_{A,t}$ works only through $Q_t^A$.
\end{proof}

\section{Bottleneck-Shock Counterfactuals}

The baseline has two external stocks, $D_t$ and $X_t$, and one endogenous domestic scarcity price.
To keep the mechanism identifiable, it is useful to compare three counterfactual economies.

\begin{definition}[External-channel counterfactuals]\label{def:counterfactuals}
For a given process for $(D_t,X_t)$, define:
\begin{enumerate}[leftmargin=1.5em]
    \item \textbf{Price-only economy:}
    \begin{equation}
        \nu_{pD}>0,\quad \nu_{pX}\ge 0,
        \qquad
        \nu_{qD}=\nu_{qX}=0.
        \label{eq:price_only}
    \end{equation}
    External states move procurement costs but do not affect physical availability.
    \item \textbf{Quantity-only economy:}
    \begin{equation}
        \nu_{pD}=\nu_{pX}=0,
        \qquad
        \nu_{qD}>0,\quad \nu_{qX}\ge 0.
        \label{eq:quantity_only}
    \end{equation}
    External states move availability and can generate an endogenous scarcity rent, but
    external procurement costs do not directly rise.
    \item \textbf{Full bottleneck economy:}
    \begin{equation}
        \nu_{pD}>0,\quad \nu_{qD}>0,
        \qquad
        \nu_{pX}\ge 0,\quad \nu_{qX}\ge 0.
        \label{eq:full_bottleneck}
    \end{equation}
    External disruption and relief affect both procurement costs and physical availability.
\end{enumerate}
\end{definition}

This decomposition is not a change in the theoretical model. It is a diagnostic exercise. The
price-only economy isolates the usual import-cost channel. The quantity-only economy isolates the
endogenous scarcity-rent channel. The full economy measures their interaction. The central
bottleneck-repair mechanism is present only to the extent that the quantity channel generates
$\chi_t>0$ and therefore changes current marginal cost and, under the local repair-benefit
restriction used above, the value of future adaptation.

These counterfactuals are useful because the same observed disruption can look like a high import
price, a low input quota, or both. The price-only case asks what happens if the shock is just more
expensive procurement. The quantity-only case asks what happens if the shock is physical scarcity.
The full case asks how the two interact.

\begin{remark}[What the counterfactuals identify]\label{rem:counterfactual_identification}
In the price-only economy, external states affect firms through the procurement price $p_t^m$ but
do not tighten imported-input capacity. If $\bar m_t$ is large enough that the cap is slack, then
$\chi_t=0$ and the model collapses to an import-cost channel without endogenous scarcity. In the
quantity-only economy, external states affect firms only by moving $\bar m_t$; any additional
private cost of imported-input use is then entirely endogenous and equals $\chi_t$. In the full
bottleneck economy, private imported-input cost decomposes into the external component $p_t^m$ and
the equilibrium scarcity component $\chi_t$.

This classification follows directly from \eqref{eq:effective_m_price} and the complementarity condition
\eqref{eq:aggregate_m_complementarity}. In the price-only case, $\nu_{qD}=\nu_{qX}=0$, so the
external states do not reduce $\bar m_t$. If the cap is slack, complementarity gives $\chi_t=0$.
In the quantity-only case, $\nu_{pD}=\nu_{pX}=0$, so $p_t^m=\bar p^m$ and movements in the
effective imported-input price beyond $\bar p^m$ must come from $\chi_t$. The full case combines
both terms in $\tilde p_t^m=p_t^m+\chi_t$.
\end{remark}

This diagnostic comparison clarifies what the model adds relative to a conventional import-price shock. If
the cap never binds, the economy has an external cost shock but no endogenous scarcity. Once the
cap binds, the model generates an internal price of access, and adaptation changes the future
importance of that internal price.

Taken together, the counterfactuals close the baseline model by identifying what is new relative to
a conventional import-price shock. A higher disruption stock $D_t$ can make the critical input more
expensive and less available, while a higher relief stock $X_t$ works in the opposite direction but
may arrive late, remain partial, and decay faster than damage. If the availability constraint
binds, the endogenous rent $\chi_t$ raises the effective cost of using the critical input. Firms
can reduce future exposure by accumulating adaptation capital, which shifts the intermediate-input
technology away from imported input use. The policy analysis can therefore be organized around one
direct current-scarcity effect and two additional conditional responses: inflation through the
Calvo marginal-cost channel and future repair through the adaptation-value channel. The current
scarcity channel is the more direct bottleneck-relief result; the inflation and repair channels
remain conditional on the local restrictions established above.

The baseline therefore establishes the minimal economy needed for the dissertation's main
mechanism. External disruption and relief determine the procurement price and physical availability
of a critical imported input. The imported-input market then determines whether an endogenous
scarcity rent appears. Calvo pricing maps effective marginal costs into inflation and price
dispersion. Firm adaptation makes recovery a dynamic private decision rather than a mechanical
return to the pre-shock state. Monetary policy enters this same equilibrium through the household
Euler equation and, when repair spending is front-loaded, through the private cost of adaptation.

The main conclusion is deliberately conditional and local. In the binding-bottleneck region, a
contractionary policy innovation can reduce current scarcity by lowering output and desired use of
the constrained input. Under the additional Calvo marginal-cost restriction, it also lowers
inflation. Under the repair-benefit restriction and on the interior repair branch, it lowers current
adaptation investment. Thus the baseline identifies the central trade-off: tighter policy can
stabilize current inflation and relieve current scarcity, but it may slow the structural reduction
of future dependence on the bottlenecked input.

The baseline also makes clear what remains outside the model at this stage. The substitute input and
the adaptation installation good are available at competitive prices, but their domestic production
is not modeled. Repair absorbs real resources through $p_{A,t}\Psi(I_t^A)$, yet the baseline does
not ask whether those resources are engineers, labor time, specialized equipment, or capital
reallocated away from current production. Firms are symmetric, and financial constraints beyond the
working-capital repair-cost wedge are not included. The next extension relaxes the most important of
these restrictions by making repair resource-based, so that recovery itself competes with current
production for domestic resources.

\section{Bridge to the Computational Chapter}

The first two parts leave the model in a form that is ready for a global numerical method. For each
policy environment $P$, the unknown object is a recursive policy function
\begin{equation}
    \mathbf y_t^P
    =
    F^P(\mathcal Z_t^P),
    \label{eq:recursive_policy_function_for_deqn}
\end{equation}
where $\mathcal Z_t^P$ is given by \eqref{eq:policy_environment_states}. The residual vector solved
in the computational chapter is built from the same equations used above:
\begin{equation}
    \mathcal R^P
    \left(
        \mathcal Z_t^P,
        F^P(\mathcal Z_t^P),
        \E_tF^P(\mathcal Z_{t+1}^P)
    \right)
    =
    \left(
        \mathcal G^B_t,
        \mathcal K^B_t,
        \boldsymbol\eta_t^B\circ\mathcal K^B_t
    \right).
    \label{eq:baseline_residual_vector}
\end{equation}
For rule-based policy, the residual vector includes the relevant Taylor rule. For discretion, it
includes the Bellman optimality conditions and the implementability constraints. For commitment, it
adds the promise-state recursion associated with the Ramsey multipliers. Thus the DEQN step does
not solve a different model; it approximates the recursive functions that make
\eqref{eq:baseline_residual_vector} small over the nonlinear state space.
Appendix~\ref{app:computational_residual_system} collects the baseline residual system explicitly:
the state vector, network output, exogenous transition block, private-equilibrium equations,
complementarity residuals, policy-rule residuals, and the additional discretion and commitment
conditions. The main text therefore keeps the economic mechanism visible, while the appendix
contains the equations needed for implementation.

The state space is nonlinear for economic reasons, not for computational decoration. External
disruption and relief arrive through persistent marked-event processes, the imported-input
constraint creates complementarity, Calvo pricing carries price dispersion as a real state, and
adaptation changes future input dependence. These objects make local perturbation around one steady
state less informative for large bottleneck episodes. The next chapter therefore uses DEQN as a
global residual-minimization method for the baseline and then for the resource-based repair and
heterogeneity extensions, with each extension adding equations to the same residual architecture.

The computational motivation is specific to this model. The relevant dynamics are not only small
fluctuations around a single steady state. The economy can move between regions in which the
imported-input cap is slack or binding, the scarcity rent is zero or positive, and repair is at a
corner or interior. These regions correspond to economically different episodes: ordinary import
cost pressure, true physical bottlenecks, and active or inactive structural repair. A local
linearization around one reference allocation cannot describe these switches well, because the main
objects of interest are precisely the state-dependent scarcity rent $\chi_t$ and the state-dependent
repair policy $I_t^A$.

The DEQN step is therefore meant to produce objects that the analytical section cannot deliver by
itself: policy functions for inflation, output, scarcity rents, repair investment, and installed
adaptation over a wide range of disruption and relief states; impulse responses after large
disruption and relief arrivals; stochastic steady states and ergodic distributions under each
policy environment; and welfare comparisons across the fixed Taylor rule, the bottleneck-adjusted
rule, discretion, and commitment. The local propositions above identify the mechanism and the
sufficient sign restrictions. The global numerical solution is needed to measure how often those
regions arise and how large the inflation--scarcity--repair trade-off is under alternative
monetary-policy environments.

\appendix

\section{Baseline Calibration}
\label{app:baseline_calibration}

This appendix records the calibration strategy for the baseline economy. The table deliberately
separates three types of discipline. First, standard New Keynesian parameters are inherited from
the persistent-supply-regime benchmark whenever the object is common to that framework and to the
present model. Second, the critical-input parameters are tied to evidence on imported-input
dependence, low input substitutability, and supply-chain pressure. Third, parameters that are not
directly observable are treated as normalizations, targets, or sensitivity parameters. They are not
presented as externally estimated numbers.

The table is written for a quarterly model. A value marked ``targeted'' means that the numerical
solution should choose the parameter to match the stated empirical moment or scenario rather than
copy a universal number from the literature. This is important for the bottleneck block: the same
model can describe semiconductors, machine parts, payment services, or logistics access, and these
objects differ in import exposure, substitutability, and relief speed.

\begingroup
\small
\setlength{\tabcolsep}{3pt}
\begin{longtable}{
>{\raggedright\arraybackslash}p{0.15\textwidth}
>{\raggedright\arraybackslash}p{0.22\textwidth}
>{\raggedright\arraybackslash}p{0.20\textwidth}
>{\raggedright\arraybackslash}p{0.35\textwidth}}
\caption{Baseline calibration map}
\label{tab:baseline_calibration}\\
\toprule
Block & Parameter & Baseline value or target & Discipline and economic interpretation \\
\midrule
\endfirsthead
\toprule
Block & Parameter & Baseline value or target & Discipline and economic interpretation \\
\midrule
\endhead
\midrule
\multicolumn{4}{r}{Continued on next page} \\
\midrule
\endfoot
\bottomrule
\endlastfoot

Frequency
&
Model period
&
Quarter
&
The baseline is quarterly, matching the monetary-policy and Calvo-pricing calibration in
Galo--Nu\~no, Renner and Scheidegger.
\\

Household
&
$\beta,\sigma,\varphi$
&
$\beta=0.9975$, $\sigma=2$, $\varphi=1$
&
Inherited from the persistent-supply-regime New Keynesian benchmark: discount factor, inverse
intertemporal elasticity, and inverse Frisch elasticity.\footnote{Galo Nu\~no, Philipp Renner and
Simon Scheidegger, ``Monetary Policy with Persistent Supply Shocks'', Banco de Espa\~na Working
Paper 2529, Table 1. The table reports $\beta=0.9975$, inverse Frisch elasticity $1$, and inverse
intertemporal elasticity $2$. \url{https://www.bde.es/f/webbe/SES/Secciones/Publicaciones/PublicacionesSeriadas/DocumentosTrabajo/25/Files/dt2529e.pdf}.}
\\

Nominal rigidity
&
$\varepsilon,\theta$
&
$\varepsilon=7$, $\theta=0.75$
&
Inherited from the same New Keynesian benchmark. $\varepsilon$ pins down the desired steady-state
markup; $\theta=0.75$ implies an average price duration of four quarters under Calvo pricing.\footnote{Galo
Nu\~no, Renner and Scheidegger, Table 1, report elasticity of substitution among varieties $7$ and
Calvo parameter $0.75$.}
\\

Productivity
&
$\bar Z,\rho_Z,\sigma_Z$
&
$\bar Z=1$, $\rho_Z=0.99$, $\sigma_Z=0.009$
&
Long-run productivity is a normalization; persistence and volatility follow the productivity block
used in the replicated NK environment.\footnote{Galo Nu\~no, Renner and Scheidegger, Table 1,
report long-run productivity $1$, TFP persistence $0.99$, and TFP innovation standard deviation
$0.009$, with shock parameters taken from Coibion et al. (2012).}
\\

Monetary rule
&
$\bar\Pi,\bar R,\phi_\Pi,\phi_Y$
&
$\bar\Pi=1$, $\bar R=\bar\Pi/\beta$, $\phi_\Pi=2$, $\phi_Y=0$ in the baseline rule; $\phi_Y=0.5$ as
sensitivity
&
Zero net inflation and inflation feedback $\phi_\Pi=2$ follow the persistent-supply-regime
benchmark. The baseline sets $\phi_Y=0$ to keep the rule close to that comparison; an output-gap
response is examined as a Taylor-rule sensitivity.\footnote{Galo Nu\~no, Renner and Scheidegger,
Table 1, use inflation target $0$ and Taylor-rule slope $2$. For the classic output-gap response,
see John B. Taylor, ``Discretion versus Policy Rules in Practice'', Carnegie-Rochester Conference
Series on Public Policy, 39, 1993, pp. 195--214, \url{https://doi.org/10.1016/0167-2231(93)90009-L}.}
\\

Production
&
$\alpha$
&
Targeted to the intermediate-input expenditure share of the sector; manufacturing benchmark:
$\alpha\simeq 0.5$
&
$\alpha$ is the cost share of the intermediate composite relative to labor in gross production. It
should be disciplined by input-output tables for the target sector rather than by the monetary
paper. The value $0.5$ is a manufacturing-style benchmark, not a universal structural constant.
\\

Critical-input exposure
&
$\omega(0)$
&
Manufacturing benchmark: $\omega(0)\simeq 0.30$
&
Initial dependence on the critical imported input is tied to imported-input reliance. A U.S.
manufacturing benchmark is that nearly one-third of intermediate inputs originate abroad, including
direct and indirect import dependence.\footnote{Andrea De Michelis and Mariano Somale, ``A Sourcing
Risk Index for U.S. Manufacturing Industries'', FEDS Notes, Board of Governors of the Federal
Reserve System, 8 September 2023. The note reports that nearly one-third of intermediate inputs
used in U.S. manufacturing originate abroad. \url{https://www.federalreserve.gov/econres/notes/feds-notes/a-sourcing-risk-index-for-u-s-manufacturing-industries-20230908.html}.}
\\

Substitution
&
$\rho$
&
$\rho=0.20$; sensitivity: $0.05$--$0.60$
&
$\rho$ is central because it determines how hard it is to replace the imported input before new
adaptation is installed. The baseline uses a low value consistent with evidence that imported and
domestic inputs can be close complements during disruption episodes.\footnote{Christoph Boehm,
Aaron Flaaen and Nitya Pandalai-Nayar, ``Input Linkages and the Transmission of Shocks:
Firm-Level Evidence from the 2011 Tohoku Earthquake'', Review of Economics and Statistics, 101(1),
2019, pp. 60--75. Their evidence points to strong complementarities and near-Leontief behavior; the
working-paper version reports an elasticity across material inputs around $0.2$.
\url{https://ideas.repec.org/a/tpr/restat/v101y2019i1p60-75.html}.}
\\

Adaptation technology
&
$\omega(A)$
&
Computational form:
$\omega(A)=\omega(0)\exp(-\kappa_AA)$
&
The analytical model requires only $\omega_A(A)<0$ and $\omega_{AA}(A)\ge0$. For computation, this
exponential form makes adaptation input-saving without giving a mechanical one-period jump to full
independence from the imported input. $\kappa_A$ is calibrated to the speed at which repair lowers
import intensity.
\\

Adaptation effectiveness
&
$\kappa_A$
&
Targeted to repair half-life; sensitivity parameter
&
There is no universal external value for the speed at which redesign, supplier switching,
recertification, and process reorganization reduce import dependence. The disciplined target is the
time needed for a firm to reduce its imported-input intensity after a disruption, using sectoral or
case-study evidence when available.
\\

Adaptation depreciation
&
$\delta_A$
&
$0.03$--$0.04$ quarterly; baseline $0.035$
&
Installed adaptation is treated like organizational and engineering know-how. A quarterly rate near
$3$--$4$ percent corresponds to annual depreciation around $12$--$15$ percent, a conventional range
for R\&D or knowledge capital.\footnote{Nadiri and Prucha estimate an annual R\&D depreciation rate
around $12$ percent for U.S. manufacturing; see BEA, ``R\&D Depreciation Rates in the 2007 R\&D
Satellite Account'', \url{https://bea.gov/sites/default/files/papers/P2007-8.pdf}. Li and Hall
argue that R\&D depreciation is often higher than the traditional $15$ percent and varies by
industry; see Wendy C. Y. Li and Bronwyn H. Hall, ``Depreciation of Business R\&D Capital'', NBER
Working Paper 22473, \url{https://www.nber.org/papers/w22473}.}
\\

Adaptation convexity
&
$\phi_A$
&
Targeted; benchmark chosen to match gradual repair, not instantaneous adjustment
&
$\phi_A$ controls how costly it is to front-load repair. It should be chosen so that, after a
large bottleneck shock, the model produces a plausible multi-quarter repair path rather than an
implausible one-period substitution away from the imported input.
\\

Repair financing
&
$\vartheta_A$
&
Baseline $0.5$; sensitivity: $0$, $0.75$, $1$
&
$\vartheta_A$ is the fraction of the repair bill that must be financed before sales revenues arrive.
The baseline uses a partial working-capital channel; the sensitivity values shut the channel down
or make repair fully pre-financed.\footnote{For micro evidence on the quantitative relevance of
working capital, see Hamilton Galindo Gil, ``Is the Working Capital Channel of Monetary Policy
Quantitatively Relevant? A Structural Estimation Approach'', Macroeconomic Dynamics, 2025. The
paper estimates a partial channel and reports that roughly three-fourths of firms' labor bill is
borrowed in the aggregate. \url{https://www.cambridge.org/core/journals/macroeconomic-dynamics/article/is-the-working-capital-channel-of-the-monetary-policy-quantitatively-relevant-a-structural-estimation-approach/BEA493ED6031DD5B74ECC90ECF3CBD86}.}
\\

Competitive supply prices
&
$p_t^d,p_{A,t}$
&
Baseline constants $p_t^d=p_{A,t}=1$; stochastic variants in robustness
&
These are normalizations in the baseline because the domestic production side of substitutes and
repair services is intentionally left outside the first model. Later resource-based repair
extensions endogenize these prices through domestic factor competition.
\\

External price normalization
&
$\bar p^m$
&
$1$
&
The steady procurement price of the imported input is a units normalization. Disruption and relief
matter through deviations generated by $\nu_{pD}D_t-\nu_{pX}X_t$.
\\

Normal input capacity
&
$\bar m$
&
$\bar m=(1+s_m)M^{d,ss}$ with $s_m=0.10$ baseline
&
The normal cap is set above unconstrained steady-state desired imported-input demand so that the
calm economy is not mechanically scarcity constrained. The slackness buffer $s_m$ is varied in
sensitivity checks because tighter normal capacity makes scarcity rents more likely.
\\

Disruption persistence
&
$\delta_D$
&
Half-life target: $24$ quarters, so $\delta_D=1-2^{-1/24}\simeq0.0285$
&
The stock of external damage is persistent. The half-life is aligned with the twenty-four-quarter
expected duration of bad times in the persistent-supply-regime benchmark, but implemented as stock
decay rather than as a discrete Markov exit probability.\footnote{Galo Nu\~no, Renner and
Scheidegger, Table 1, use transition probability $p_{21}=1/24$ from negative supply times to
normal times.}
\\

Disruption arrival rate
&
$\bar\lambda_D$
&
$1/48$
&
The baseline arrival rate gives one major disruption episode every forty-eight quarters in the
normal external environment. This imports the frequency scale of the normal-to-bad transition in
the persistent-regime benchmark while replacing fixed regimes with marked events.\footnote{Galo
Nu\~no, Renner and Scheidegger, Table 1, use transition probability $p_{12}=1/48$ from normal to
negative supply times.}
\\

Relief persistence
&
$\delta_X$
&
Half-life target: $12$ quarters, so $\delta_X=1-2^{-1/12}\simeq0.056$
&
Relief is made less durable than damage. This encodes the idea that workaround routes, settlement
channels, and alternative suppliers may be partial and fragile. The exact half-life is a scenario
discipline, not an externally estimated structural constant.
\\

Relief arrival rate
&
$\bar\lambda_X,\beta_X$
&
$\bar\lambda_X$ targeted; $\beta_X>0$
&
Relief should be delayed and state dependent. $\beta_X>0$ makes relief arrivals more likely after
damage accumulates, while $\bar\lambda_X$ controls the unconditional background frequency. These
parameters should be chosen to match the observed speed at which external workarounds appear after
large disruptions.
\\

Event-intensity persistence
&
$\rho_{\lambda D},\rho_{\lambda X}$
&
Sensitivity range: $0$--$0.90$
&
Constant intensities correspond to $0$. Positive persistence allows clusters of disruption or
relief events. The empirical discipline should come from supply-chain pressure indicators, because
global bottleneck pressure is visibly persistent and episode-based rather than iid.\footnote{Gianluca
Benigno, Julian di Giovanni, Jan J. J. Groen and Adam I. Noble, ``The GSCPI: A New Barometer of
Global Supply Chain Pressures'', Federal Reserve Bank of New York Staff Report 1017, 2022.
\url{https://www.newyorkfed.org/medialibrary/media/research/staff_reports/sr1017.pdf}.}
\\

Event-size distributions
&
$m^D,\zeta^X$
&
Targeted marked distributions
&
The marks determine how large individual disruption and relief events are. They should be
calibrated from event studies or supply-chain-pressure episodes: adverse marks should be rare and
sometimes large; relief marks should be smaller and more gradual.
\\

Mapping to procurement price
&
$\nu_{pD},\nu_{pX}$
&
Targeted to input-price and shipping-cost responses
&
These parameters map external damage and relief into the external procurement price $p_t^m$.
They should be disciplined by sectoral import prices, freight rates, or supply-chain cost indices.
They are not generic NK parameters.
\\

Mapping to physical availability
&
$\nu_{qD},\nu_{qX}$
&
Targeted to quantity shortfalls; impose $\nu_{qD}\ge\nu_{pD}$ and $\nu_{qX}>\nu_{pX}$
&
These parameters determine how strongly disruption and relief move the physical input cap
$\bar m_t$. The inequalities encode the baseline bottleneck hypothesis: disruption primarily
tightens availability, and relief restores access more than it reduces procurement cost.
\\

Shock volatilities
&
$\sigma_{\lambda D},\sigma_{\lambda X}$
&
Estimated or sensitivity parameters
&
These govern volatility in event intensities, not ordinary one-period cost-push noise. They should
be disciplined by the volatility of supply-chain pressure indicators or by scenario analysis around
large disruption episodes.
\\

Initial state
&
$D_0,X_0,A_0,\Delta_{-1}$
&
Steady state: $D_0=X_0=A_0=0$, $\Delta_{-1}=1$; shock experiments move $D_0$ or event arrivals
&
The calm starting point has no accumulated disruption, no accumulated external relief, no installed
repair stock, and no price dispersion. Large-shock experiments then introduce disruption events and
trace how scarcity and repair evolve.
\\

\end{longtable}
\endgroup

The calibration is intentionally modular. The NK backbone is fixed by the replicated
persistent-supply-regime environment. The bottleneck block is disciplined by imported-input
exposure, low substitutability, and supply-chain pressure evidence. The repair block is disciplined
by the speed and cost of reducing input dependence, with depreciation tied to knowledge-capital
evidence and financing tied to working-capital evidence. The remaining parameters are reported as
targets or sensitivity ranges because assigning a single number without sectoral evidence would
create false precision. In the quantitative chapter, the baseline should therefore report not only
one benchmark calibration, but also sensitivity to $\rho$, $\bar m$, $\delta_D$, $\delta_X$,
$\phi_A$, and $\vartheta_A$, since these parameters directly control the strength of the
inflation--scarcity--repair trade-off.

\subsection{Sensitivity Design and Mechanism Diagnostics}
\label{app:sensitivity_design}

The quantitative section should not treat sensitivity analysis as a mechanical robustness appendix.
In this model, changing parameters changes the economic type of bottleneck episode. The purpose of
the sensitivity exercises is therefore to identify which force drives the policy trade-off:
external crisis persistence, the physical tightness of the input constraint, the speed of external
relief, or the ability of firms to repair their own technology.

\begingroup
\small
\setlength{\tabcolsep}{3pt}
\begin{longtable}{
>{\raggedright\arraybackslash}p{0.20\textwidth}
>{\raggedright\arraybackslash}p{0.23\textwidth}
>{\raggedright\arraybackslash}p{0.23\textwidth}
>{\raggedright\arraybackslash}p{0.26\textwidth}}
\caption{Sensitivity experiments for the baseline quantitative model}
\label{tab:sensitivity_design}\\
\toprule
Experiment & Parameters moved & Economic question & Expected diagnostic pattern \\
\midrule
\endfirsthead
\toprule
Experiment & Parameters moved & Economic question & Expected diagnostic pattern \\
\midrule
\endhead
\midrule
\multicolumn{4}{r}{Continued on next page} \\
\midrule
\endfoot
\bottomrule
\endlastfoot

Crisis depth
&
Larger disruption marks $m^D$; higher $\nu_{qD}$ and $\nu_{pD}$
&
What changes when the same kind of event becomes more severe?
&
Inflation, scarcity rents, and output losses rise together. The distinctive test is whether repair
investment rises because repair becomes more valuable, or falls because the crisis destroys current
cash-flow and demand conditions.
\\

Crisis persistence
&
Lower $\delta_D$; higher $\rho_{\lambda D}$; higher $\kappa_D^\lambda$
&
Does the policy trade-off become stronger when bad external conditions last longer or cluster?
&
Longer expected bottleneck duration raises the continuation value $Q_t^A$ of adaptation. Tight
policy then has a larger intertemporal cost if it depresses repair exactly when repair has become
more valuable.
\\

Pure price pressure
&
High $\nu_{pD}$, low $\nu_{qD}$
&
How much of the inflation response is just an external import-cost shock?
&
The scarcity rent $\chi_t$ should play a smaller role. If the policy trade-off weakens here, the
central mechanism is not merely imported inflation but the binding physical bottleneck.
\\

Pure quantity bottleneck
&
High $\nu_{qD}$, low $\nu_{pD}$
&
What happens when the crisis mainly cuts availability rather than raising procurement prices?
&
The endogenous rent $\chi_t$ becomes the key state. Monetary tightening can lower current scarcity
by reducing input demand, but this may be precisely the case in which delayed repair is most costly.
\\

Relief speed
&
Higher $\bar\lambda_X$, higher $\beta_X$, larger relief marks $\zeta^X$
&
Does faster external normalization reduce the need for firm-level repair?
&
Faster relief lowers expected future scarcity and should reduce $Q_t^A$. The repair channel of
monetary policy becomes weaker because firms expect the outside environment to heal without as much
internal adaptation.
\\

Relief fragility
&
Higher $\delta_X$; lower $\zeta^X$; lower $\nu_{qX}$
&
What changes when workarounds are temporary, partial, or unreliable?
&
External relief improves current conditions but does not eliminate future bottleneck risk. This
can produce recurrent scarcity episodes and a stronger value of maintaining adaptation even after
the first crisis wave fades.
\\

Normal capacity slack
&
Lower or higher $s_m$ in $\bar m=(1+s_m)M^{d,ss}$
&
How close is the normal economy to the physical capacity constraint?
&
With large slack, crises must be deep before $\chi_t>0$. With small slack, even moderate disruption
creates scarcity rents. This separates an ordinary cost-push episode from a true bottleneck episode.
\\

Input substitutability
&
Lower or higher $\rho$
&
How much does technological substitutability discipline the strength of the bottleneck?
&
Low $\rho$ makes the imported input hard to replace within the current technology. Scarcity rents
and repair values become more sensitive to shocks. High $\rho$ weakens the bottleneck and makes
the model closer to a standard relative-price shock.
\\

Initial import exposure
&
Higher or lower $\omega(0)$
&
Does the trade-off depend on how import-dependent the economy is before the crisis?
&
Higher initial exposure raises the level of imported-input demand and makes the cap bind more
often. It also increases the payoff from adaptation, so the repair margin should become more
important.
\\

Adaptation effectiveness
&
Higher or lower $\kappa_A$
&
Does repair actually release the economy from the bottleneck?
&
When $\kappa_A$ is high, a unit of repair sharply lowers future import dependence, so policy-induced
repair delays have large future costs. When $\kappa_A$ is low, repair is costly but not very
effective, and the welfare value of preserving repair is smaller.
\\

Repair speed and convexity
&
Higher or lower $\phi_A$; alternative repair half-lives through $\delta_A$
&
Is recovery limited by installation costs or by depreciation of installed capabilities?
&
High $\phi_A$ makes repair gradual and creates stronger intertemporal smoothing. High $\delta_A$
makes repair less durable and weakens the value of past adaptation. These parameters determine
whether the economy escapes the bottleneck or repeatedly falls back into dependence.
\\

Repair financing exposure
&
Higher or lower $\vartheta_A$
&
How much of the policy trade-off comes from the direct financing cost of repair?
&
When $\vartheta_A=0$, monetary policy affects repair only through continuation values and demand.
When $\vartheta_A$ is high, the same rate increase directly raises the private cost of repair, so
the repair slowdown channel is stronger.
\\

Policy feedback strength
&
Higher or lower $\phi_\Pi$ and $\phi_Y$
&
Does stronger nominal stabilization deepen the repair trade-off?
&
A more aggressive inflation response should reduce current inflation and scarcity pressure faster,
but may also compress demand and repair incentives more strongly. The key output is not only lower
inflation volatility, but the path of $I_t^A$ and $A_{t+1}$.
\\

\end{longtable}
\endgroup

These experiments should be reported around four objects: inflation $\Pi_t$, the scarcity rent
$\chi_t$, repair investment $I_t^A$, and installed adaptation $A_t$. Output and consumption are
necessary for welfare, but the model's new mechanism is identified by the joint movement of
$\chi_t$ and $I_t^A$. A parameter change is economically important when it changes the sign,
timing, or persistence of this pair. For example, a deep but short price shock can raise inflation
without generating a large repair motive, while a persistent quantity bottleneck can generate a
smaller immediate inflation spike but a much larger value of structural adaptation.

The most informative patterns are therefore not single impulse responses but comparisons across
policy environments. A fixed Taylor rule can stabilize ordinary demand pressure while ignoring
bottleneck-state dependence. A bottleneck-adjusted rule corrects the natural-rate intercept but
still does not optimize repair incentives. Discretion internalizes current implementability
conditions but cannot commit to future repair-supporting policy. Commitment can influence current
repair values through expectations. Sensitivity analysis should ask when these four environments
rank differently, because those reversals reveal whether the main issue is nominal stabilization,
current bottleneck relief, or delayed structural escape.

\section{Baseline Computational Residuals}
\label{app:computational_residual_system}

This appendix rewrites the baseline economy as the residual system used by the global numerical
solution. It does not introduce a new model. The purpose is purely computational: the equations
derived in the main text are collected in the form
\begin{equation}
    \mathcal R^P
    \left(
        \mathcal Z_t^P,
        F_\nu^P(\mathcal Z_t^P),
        \E_t F_\nu^P(\mathcal Z_{t+1}^P)
    \right)
    =
    0,
    \label{eq:app_deqn_master_residual}
\end{equation}
where $P$ indexes the policy environment and $F_\nu^P$ is the neural-network approximation to the
recursive policy function. In the rule-based cases, the residual vector is the private-equilibrium
system plus a policy rule. Under discretion and commitment, it also contains the corresponding
optimal-policy first-order conditions.

\subsection{State Vector and Network Output}

For the fixed Taylor rule, the bottleneck-adjusted Taylor rule, and discretion, the physical
recursive state is
\begin{equation}
    \mathcal Z_t
    =
    (D_t,X_t,\lambda_t^D,\lambda_t^X,Z_t,p_t^d,p_{A,t},A_t,\Delta_{t-1}).
    \label{eq:app_deqn_state}
\end{equation}
For commitment, the recursive implementation augments the same physical state with promise
variables. The computational promise variables are scaled promises, not raw lagged multipliers:
\begin{equation}
    \mathcal Z_t^C
    =
    (\mathcal Z_t,\mathbf p_t^{E},\mathbf p_t^{P},\mathbf p_t^{A}).
    \label{eq:app_deqn_commitment_state}
\end{equation}
The exact dimension of the promise block depends on the chosen Ramsey implementation, but the
economic content is the same as in the main text: commitment makes lagged promises attached to
forward-looking private constraints part of the state. Because the promises are scaled by the
current normalizers when they are created, lagged consumption is not a separate state variable in
this compact implementation.

For computation it is convenient to let the network return both allocation variables and auxiliary
objects that could otherwise be substituted out:
\begin{equation}
    F_\nu^P(\mathcal Z_t^P)
    =
    \mathbf y_t^P
    =
    (C_t,Y_t,N_t,w_t,\Lambda_t,\mc_t,p_t^x,\mathcal X_t,M_t,S_t,
    \chi_t,\Pi_t,p_t^*,\Delta_t,I_t^A,A_{t+1},
    \mathcal S_t^p,\mathcal F_t^p,Q_t^A,R_t).
    \label{eq:app_deqn_output_vector}
\end{equation}
This output vector is not unique. A code implementation may eliminate auxiliary variables such as
$\Lambda_t$, $p_t^x$, $\mc_t$, or $\mathcal X_t$ by substitution. The residual system below keeps
them explicit because doing so makes each economic restriction visible and easier to test.

\subsection{Exogenous Transition and Input-Condition Block}

The exogenous transition block generates next-period states used in the expectation operator:
\begin{align}
    D_{t+1}
    &=
    (1-\delta_D)D_t
    +
    \sum_{j=1}^{N_{t+1}^D}m_{t+1,j}^D,
    &
    N_{t+1}^D\mid\lambda_t^D
    &\sim \mathrm{Poisson}(\lambda_t^D),
    \label{eq:app_deqn_D_transition}\\
    X_{t+1}
    &=
    (1-\delta_X)X_t
    +
    \sum_{j=1}^{N_{t+1}^X}\zeta_{t+1,j}^X,
    &
    N_{t+1}^X\mid\lambda_t^X
    &\sim \mathrm{Poisson}(\lambda_t^X).
    \label{eq:app_deqn_X_transition}
\end{align}
When event intensities are stochastic, their laws are
\begin{align}
    \log\lambda_{t+1}^D
    &=
    (1-\rho_{\lambda D})\log\bar\lambda_D
    +
    \rho_{\lambda D}\log\lambda_t^D
    +
    \kappa_D^\lambda D_t
    +
    \sigma_{\lambda D}\varepsilon_{t+1}^{\lambda D},
    \label{eq:app_deqn_lambda_D}\\
    \log\lambda_{t+1}^X
    &=
    (1-\rho_{\lambda X})\log\bar\lambda_X
    +
    \rho_{\lambda X}\log\lambda_t^X
    +
    \beta_XD_t
    +
    \sigma_{\lambda X}\varepsilon_{t+1}^{\lambda X}.
    \label{eq:app_deqn_lambda_X}
\end{align}
The same simulated or quadrature-based transition also draws $Z_{t+1}$, $p_{t+1}^d$, and
$p_{A,t+1}$ from their Markov laws. Given $(D_t,X_t)$, the current imported-input procurement
conditions are computed as
\begin{align}
    p_t^m
    &=
    \bar p^m
    \exp(\nu_{pD}D_t-\nu_{pX}X_t),
    \label{eq:app_deqn_pm}\\
    \bar m_t
    &=
    \bar m
    \exp(-\nu_{qD}D_t+\nu_{qX}X_t).
    \label{eq:app_deqn_mbar}
\end{align}
Expectations in the residuals below are evaluated over this transition:
\begin{equation}
    \E_t h(\mathcal Z_{t+1}^P)
    =
    \int
    h\!\left(T^P(\mathcal Z_t^P,\varepsilon_{t+1})\right)
    \dd \mathbb P(\varepsilon_{t+1}\mid\mathcal Z_t^P),
    \label{eq:app_deqn_expectation_operator}
\end{equation}
where $T^P$ includes both exogenous laws and endogenous state updates, such as
$A_{t+1}$ and $\Delta_t$.

\subsection{Private-Equilibrium Residuals}

The household residuals are
\begin{align}
    \mathcal R_t^\Lambda
    &\equiv
    \Lambda_t-C_t^{-\sigma},
    \label{eq:app_deqn_R_lambda}\\
    \mathcal R_t^N
    &\equiv
    N_t^\varphi-w_t\Lambda_t,
    \label{eq:app_deqn_R_labor_supply}\\
    \mathcal R_t^E
    &\equiv
    1-\beta R_t
    \E_t
    \left[
        \frac{\Lambda_{t+1}}{\Lambda_t}
        \frac{1}{\Pi_{t+1}}
    \right].
    \label{eq:app_deqn_R_euler}
\end{align}
The imported-input and production-cost residuals are
\begin{align}
    \mathcal R_t^{\tilde m}
    &\equiv
    \tilde p_t^m-(p_t^m+\chi_t),
    \label{eq:app_deqn_R_effective_price}\\
    \mathcal R_t^{x,p}
    &\equiv
    p_t^x
    -
    \left[
        \omega(A_t)^\rho(\tilde p_t^m)^{1-\rho}
        +
        (1-\omega(A_t))^\rho(p_t^d)^{1-\rho}
    \right]^{\frac{1}{1-\rho}},
    \label{eq:app_deqn_R_px}\\
    \mathcal R_t^{mc}
    &\equiv
    \mc_t
    -
    \frac{1}{Z_t}
    \left(\frac{w_t}{1-\alpha}\right)^{1-\alpha}
    \left(\frac{p_t^x}{\alpha}\right)^\alpha .
    \label{eq:app_deqn_R_mc}
\end{align}
Given $p_t^x$ and $\mc_t$, aggregate input requirements satisfy
\begin{align}
    \mathcal R_t^{X}
    &\equiv
    \mathcal X_t
    -
    \alpha\frac{\mc_t\Delta_tY_t}{p_t^x},
    \label{eq:app_deqn_R_X}\\
    \mathcal R_t^{M}
    &\equiv
    M_t
    -
    \mathcal X_t
    \omega(A_t)^\rho
    \left(\frac{p_t^x}{\tilde p_t^m}\right)^\rho,
    \label{eq:app_deqn_R_M}\\
    \mathcal R_t^{S}
    &\equiv
    S_t
    -
    \mathcal X_t
    (1-\omega(A_t))^\rho
    \left(\frac{p_t^x}{p_t^d}\right)^\rho,
    \label{eq:app_deqn_R_S}\\
    \mathcal R_t^{N,d}
    &\equiv
    N_t
    -
    (1-\alpha)\frac{\mc_t\Delta_tY_t}{w_t}.
    \label{eq:app_deqn_R_N_demand}
\end{align}
The Calvo residuals are
\begin{align}
    \mathcal R_t^{idx}
    &\equiv
    1-(1-\theta)(p_t^*)^{1-\varepsilon}-\theta\Pi_t^{\varepsilon-1},
    \label{eq:app_deqn_R_calvo_index}\\
    \mathcal R_t^{\Delta}
    &\equiv
    \Delta_t-(1-\theta)(p_t^*)^{-\varepsilon}
    -\theta\Pi_t^\varepsilon\Delta_{t-1},
    \label{eq:app_deqn_R_delta}\\
    \mathcal R_t^{S,p}
    &\equiv
    \mathcal S_t^p-\mc_tY_t
    -
    \theta\E_tM_{t,t+1}\Pi_{t+1}^{\varepsilon}\mathcal S_{t+1}^p,
    \label{eq:app_deqn_R_Sp}\\
    \mathcal R_t^{F,p}
    &\equiv
    \mathcal F_t^p-Y_t
    -
    \theta\E_tM_{t,t+1}\Pi_{t+1}^{\varepsilon-1}\mathcal F_{t+1}^p,
    \label{eq:app_deqn_R_Fp}\\
    \mathcal R_t^{*}
    &\equiv
    p_t^*
    -
    \frac{\varepsilon}{\varepsilon-1}
    \frac{\mathcal S_t^p}{\mathcal F_t^p}.
    \label{eq:app_deqn_R_reset}
\end{align}
The adaptation and aggregate-resource residuals are
\begin{align}
    \mathcal R_t^{A,L}
    &\equiv
    A_{t+1}-(1-\delta_A)A_t-I_t^A,
    \label{eq:app_deqn_R_A_law}\\
    \mathcal R_t^Q
    &\equiv
    Q_t^A
    -
    \E_tM_{t,t+1}
    \left[
        -\C_{A,t+1}^{\Delta}
        +(1-\delta_A)Q_{t+1}^A
    \right],
    \label{eq:app_deqn_R_Q}\\
    \mathcal R_t^Y
    &\equiv
    Y_t-C_t-p_t^mM_t-p_t^dS_t-p_{A,t}\Psi(I_t^A).
    \label{eq:app_deqn_R_resource}
\end{align}
The real discount factor used in the forward-looking residuals is
\begin{equation}
    M_{t,t+1}=\beta\frac{\Lambda_{t+1}}{\Lambda_t}.
    \label{eq:app_deqn_sdf}
\end{equation}

\subsection{Complementarity Residuals}

The two occasionally binding restrictions are the imported-input cap and the non-negativity
condition for repair:
\begin{align}
    0\le \chi_t
    &\perp
    \bar m_t-M_t\ge0,
    \label{eq:app_deqn_m_mcp}\\
    0\le I_t^A
    &\perp
    \Omega_t^Ap_{A,t}\Psi'(I_t^A)-Q_t^A\ge0,
    \qquad
    \Omega_t^A=(1-\vartheta_A)+\vartheta_AR_t.
    \label{eq:app_deqn_I_mcp}
\end{align}
One implementation can keep these as mixed-complementarity restrictions. Another can use a
Fischer--Burmeister residual
\begin{equation}
    \Phi^{FB}(a,b)=\sqrt{a^2+b^2}-a-b,
    \label{eq:app_deqn_FB}
\end{equation}
so that the two nonsmooth residuals are
\begin{align}
    \mathcal R_t^{MCP,m}
    &\equiv
    \Phi^{FB}(\chi_t,\bar m_t-M_t),
    \label{eq:app_deqn_R_m_mcp}\\
    \mathcal R_t^{MCP,A}
    &\equiv
    \Phi^{FB}(I_t^A,\Omega_t^Ap_{A,t}\Psi'(I_t^A)-Q_t^A).
    \label{eq:app_deqn_R_A_mcp}
\end{align}
This is the computational representation of the model's economically important switching regions:
slack versus binding imported-input capacity and inactive versus active repair.

\subsection{Rule-Based Policy Residuals}

For the fixed-intercept Taylor rule, the policy residual is
\begin{equation}
    \mathcal R_t^{F,rule}
    \equiv
    \log R_t
    -
    \left[
        \log\bar R
        +
        \phi_\Pi\log\left(\frac{\Pi_t}{\bar\Pi}\right)
        +
        \phi_Y\log\left(\frac{Y_t}{Y_t^n}\right)
        +
        \varepsilon_t^R
    \right].
    \label{eq:app_deqn_fixed_rule}
\end{equation}
For the bottleneck-adjusted Taylor rule,
\begin{equation}
    \mathcal R_t^{BA,rule}
    \equiv
    \log R_t
    -
    \left[
        \log R_t^n
        +
        \phi_\Pi\log\left(\frac{\Pi_t}{\bar\Pi}\right)
        +
        \phi_Y\log\left(\frac{Y_t}{Y_t^n}\right)
        +
        \varepsilon_t^R
    \right],
    \label{eq:app_deqn_ba_rule}
\end{equation}
where $Y_t^n$ and $R_t^n$ are computed from the flexible-price benchmark at the same physical state.
Let the smooth private-equilibrium residuals be collected by block:
\begin{align}
    \mathcal R_t^{hh}
    &=
    (\mathcal R_t^\Lambda,\mathcal R_t^N,\mathcal R_t^E),
    \label{eq:app_deqn_hh_residual_vector}\\
    \mathcal R_t^{cost}
    &=
    (\mathcal R_t^{\tilde m},\mathcal R_t^{x,p},\mathcal R_t^{mc}),
    \label{eq:app_deqn_cost_residual_vector}\\
    \mathcal R_t^{input}
    &=
    (\mathcal R_t^X,\mathcal R_t^M,\mathcal R_t^S,\mathcal R_t^{N,d}),
    \label{eq:app_deqn_input_residual_vector}\\
    \mathcal R_t^{calvo}
    &=
    (\mathcal R_t^{idx},\mathcal R_t^\Delta,\mathcal R_t^{S,p},
    \mathcal R_t^{F,p},\mathcal R_t^*),
    \label{eq:app_deqn_calvo_residual_vector}\\
    \mathcal R_t^{adapt}
    &=
    (\mathcal R_t^{A,L},\mathcal R_t^Q,\mathcal R_t^Y).
    \label{eq:app_deqn_adapt_residual_vector}
\end{align}
The private residual vector is
\begin{equation}
    \mathcal R_t^{priv}
    =
    (\mathcal R_t^{hh},\mathcal R_t^{cost},\mathcal R_t^{input},
    \mathcal R_t^{calvo},\mathcal R_t^{adapt}).
    \label{eq:app_deqn_private_residual_vector}
\end{equation}
The full rule-based residual vector is therefore
\begin{equation}
    \mathcal R_t^{P}
    =
    \left(
        \mathcal R_t^{priv},
        \mathcal R_t^{MCP,m},
        \mathcal R_t^{MCP,A},
        \mathcal R_t^{P,rule}
    \right),
    P\in\{F,BA\}.
    \label{eq:app_deqn_rule_residual_vector}
\end{equation}

\subsection{Discretion and Commitment Residuals}

The rule-based systems close the private economy with one policy rule. Optimal policy instead
replaces the rule residual by first-order conditions of the policymaker's implementable-equilibrium
problem. To state these equations without hiding the occasionally binding constraints, define a
branch $b$ as one element of
\[
    b\in
    \{
    \text{binding cap},\text{slack cap}
    \}
    \times
    \{
    \text{interior repair},\text{repair corner}
    \}.
\]
For each branch, let $\mathcal H_t^b$ collect the smooth private residuals and the active
equalities from that branch:
\begin{equation}
    \mathcal H_t^b
    \equiv
    \left(
        \mathcal R_t^{priv},
        \mathcal E_t^{M,b},
        \mathcal E_t^{A,b}
    \right),
    \label{eq:app_opt_branch_residual}
\end{equation}
where
\begin{align}
    \mathcal E_t^{M,b}
    &=
    \begin{cases}
    \bar m_t-M_t, & \text{binding imported-input cap},\\
    \chi_t, & \text{slack imported-input cap},
    \end{cases}
    \label{eq:app_opt_cap_branch_eq}\\
    \mathcal E_t^{A,b}
    &=
    \begin{cases}
    \Omega_t^Ap_{A,t}\Psi'(I_t^A)-Q_t^A, & \text{interior repair},\\
    I_t^A, & \text{repair corner}.
    \end{cases}
    \label{eq:app_opt_repair_branch_eq}
\end{align}
The inactive inequalities are retained as branch-admissibility conditions:
\begin{align}
    \chi_t>0,\ \bar m_t-M_t=0
    &\quad\text{on the binding-cap branch}, \nonumber\\
    \chi_t=0,\ \bar m_t-M_t>0
    &\quad\text{on the slack-cap branch}, \nonumber\\
    I_t^A>0,\ \Omega_t^Ap_{A,t}\Psi'(I_t^A)-Q_t^A=0
    &\quad\text{on the interior-repair branch}, \nonumber\\
    I_t^A=0,\ \Omega_t^Ap_{A,t}\Psi'(0)-Q_t^A\ge0
    &\quad\text{on the repair-corner branch}.
    \label{eq:app_opt_branch_admissibility}
\end{align}
Thus the optimal-policy residuals below are written for a fixed active branch. At a switching
boundary the mixed-complementarity residuals \eqref{eq:app_deqn_R_m_mcp}--\eqref{eq:app_deqn_R_A_mcp}
replace the smooth branch equation.

Let $\mathbf a_t$ denote the current implementable allocation vector, equal to the private
equilibrium vector in \eqref{eq:policy_choice_vector}. The current physical state evolves according
to
\begin{equation}
    \mathcal Z_{t+1}
    =
    \mathcal T(\mathcal Z_t,\mathbf a_t,\varepsilon_{t+1}),
    \label{eq:app_opt_state_transition}
\end{equation}
where the exogenous components follow Assumptions~\ref{ass:D_law}--\ref{ass:admissible_env}, while
the endogenous predetermined components are $A_{t+1}$ and $\Delta_t$. This notation makes clear
which current choices become tomorrow's state variables.

\subsubsection*{Discretion}

Under discretion, the central bank takes future discretionary policy functions as given. For a
branch $b$, the unknown recursive object is
\begin{equation}
    F_\nu^{D,b}(\mathcal Z_t)
    =
    \left(
        \mathbf a_t^D,
        V_t^D,
        \boldsymbol\mu_t^{D,b}
    \right),
    \label{eq:app_discretion_network_output}
\end{equation}
where $\boldsymbol\mu_t^{D,b}$ is the multiplier vector on $\mathcal H_t^b=0$. The discretionary
Bellman residual is
\begin{equation}
    \mathcal R_t^{D,V}
    \equiv
    V_t^D
    -
    U(C_t,N_t)
    -
    \beta
    \E_t
    V_\nu^D(\mathcal Z_{t+1})
    =
    0.
    \label{eq:app_deqn_discretion_bellman}
\end{equation}
The stationarity residual for the implementable allocation is the vector
\begin{equation}
    \mathcal R_t^{D,a}
    \equiv
    U_{\mathbf a,t}
    +
    \left(\mathcal H_{\mathbf a,t}^{b}\right)'
    \boldsymbol\mu_t^{D,b}
    +
    \beta
    \E_t
    \left[
        V_{\mathcal Z}^D(\mathcal Z_{t+1})
        \mathcal T_{\mathbf a,t}
    \right]
    =
    0.
    \label{eq:app_deqn_discretion_stationarity}
\end{equation}
Here $U_{\mathbf a,t}$ is zero for allocation components that do not enter household utility
directly, and $\mathcal T_{\mathbf a,t}$ is nonzero only for choices that affect the next
predetermined state, mainly $A_{t+1}$ and $\Delta_t$. Equation
\eqref{eq:app_deqn_discretion_stationarity} is the complete discretionary FOC system in vector
form. The scalar margins for $R_t$, $I_t^A$, $\Pi_t$, and $\Delta_t$ are expanded in
Appendix~\ref{app:expanded_optimal_policy_focs}.

The discretionary envelope residual is
\begin{equation}
    \mathcal R_t^{D,Z}
    \equiv
    V_{\mathcal Z}^D(\mathcal Z_t)
    -
    \left[
        \left(\mathcal H_{\mathcal Z,t}^{b}\right)'
        \boldsymbol\mu_t^{D,b}
        +
        \beta
        \E_t
        \left[
            V_{\mathcal Z}^D(\mathcal Z_{t+1})
            \mathcal T_{\mathcal Z,t}
        \right]
    \right]
    =
    0.
    \label{eq:app_deqn_discretion_envelope}
\end{equation}
This envelope equation is useful for describing the continuous-time logic of the recursive problem,
but the baseline computational implementation does not add it as a separate loss block. Instead,
$V_{\mathcal Z}^D$ is obtained by automatic differentiation of the same differentiable value
network that appears in \eqref{eq:app_deqn_discretion_bellman} and
\eqref{eq:app_deqn_discretion_stationarity}. If one used a separate value-gradient network, then
\eqref{eq:app_deqn_discretion_envelope} would be the residual that pins that network down.

The full discretionary residual vector used in the baseline DEQN implementation on branch $b$ is
therefore
\begin{equation}
    \mathcal R_t^{D,b}
    =
    \left(
        \mathcal H_t^b,
        \mathcal R_t^{D,V},
        \mathcal R_t^{D,a}
    \right),
    \label{eq:app_deqn_discretion_full_residual}
\end{equation}
together with the branch-admissibility inequalities in \eqref{eq:app_opt_branch_admissibility}.
This is the discretion analogue of the rule-based residual vector
\eqref{eq:app_deqn_rule_residual_vector}: the policy rule is replaced by Bellman optimality and
stationarity. The envelope condition \eqref{eq:app_deqn_discretion_envelope} is retained as an
implementation diagnostic and as the equation that would be imposed directly in a formulation with a
separate value-gradient network.

\subsubsection*{Commitment}

Under commitment, the central bank chooses an intertemporal implementable plan. Current choices
enter not only current constraints, but also earlier forward-looking constraints through promises.
For a branch $b$, define the Ramsey multiplier vector $\boldsymbol\mu_t^{C,b}$ on
$\mathcal H_t^b=0$. The one-period promise object carried into period $t+1$ is
\begin{equation}
    \mathbf p_{t+1}^C
    =
    \mathcal P(\boldsymbol\mu_t^{C,b},\mathbf a_t),
    \label{eq:app_commitment_promise_map}
\end{equation}
where $\mathcal P(\cdot)$ selects the multipliers attached to forward-looking private constraints
and applies the current scaling implied by the normalization of those constraints. In particular,
the promise carried from period $t$ to $t+1$ absorbs the date-$t$ inverse marginal-utility factor
that would otherwise require carrying lagged consumption as a separate auxiliary state. The
household Euler equation, the Calvo pricing recursions, and the adaptation-value recursion each
contribute one or more scaled promise components.
The recursive commitment state is
\begin{equation}
    \mathcal Z_t^C
    =
    (\mathcal Z_t,\mathbf p_t^C),
    \qquad
    \mathcal Z_{t+1}^C
    =
    \left(
        \mathcal T(\mathcal Z_t,\mathbf a_t,\varepsilon_{t+1}),
        \mathcal P(\boldsymbol\mu_t^{C,b},\mathbf a_t)
    \right).
    \label{eq:app_deqn_promise_transition}
\end{equation}
Thus $C_t$ remains a current control. What enters the next state is the scaled promise object
constructed from current multipliers and current allocation normalizers.
A commitment network therefore returns
\begin{equation}
    F_\nu^{C,b}(\mathcal Z_t^C)
    =
    \left(
        \mathbf a_t^C,
        \boldsymbol\mu_t^{C,b},
        \mathbf p_{t+1}^C
    \right).
    \label{eq:app_commitment_network_output}
\end{equation}
The promise-transition residual is
\begin{equation}
    \mathcal R_t^{C,p}
    \equiv
    \mathbf p_{t+1}^C
    -
    \mathcal P(\boldsymbol\mu_t^{C,b},\mathbf a_t)
    =
    0.
    \label{eq:app_commitment_promise_residual}
\end{equation}

The Ramsey stationarity residual for the current implementable allocation is
\begin{equation}
    \mathcal R_t^{C,a}
    \equiv
    U_{\mathbf a,t}
    +
    \left(\mathcal H_{\mathbf a,t}^{b}\right)'
    \boldsymbol\mu_t^{C,b}
    +
    \boldsymbol\varrho_{t-1}^{C}(\mathbf a_t)
    +
    \beta
    \E_t
    \left[
        \left(\mathcal H_{\mathcal Z,t+1}^{b}\right)'
        \boldsymbol\mu_{t+1}^{C,b}
        \mathcal T_{\mathbf a,t}
    \right]
    =
    0.
    \label{eq:app_deqn_commitment_stationarity}
\end{equation}
The term $\boldsymbol\varrho_{t-1}^{C}(\mathbf a_t)$ is the derivative of the previous-period
forward-looking constraints with respect to the current allocation, written after applying the same
scaling normalization used in \eqref{eq:app_commitment_promise_map}:
\begin{equation}
    \boldsymbol\varrho_{t-1}^{C}(\mathbf a_t)
    \equiv
    \left(
        \mathbf p_t^C
    \right)'
    \frac{\partial}{\partial\mathbf a_t}
    \widetilde{\mathcal H}_{t-1}^{b}.
    \label{eq:app_commitment_lagged_promise_term}
\end{equation}
Here $\widetilde{\mathcal H}_{t-1}^{b}$ denotes the forward-looking part of the previous-period
constraint system expressed in the scaled variables. This notation is what keeps the recursive
state compact: raw lagged multipliers and lagged consumption are not carried separately because
their product is summarized by $\mathbf p_t^C$.
This is the exact object that makes commitment different from discretion. Under discretion, the
central bank takes future policy functions as given and has no inherited promise term. Under
commitment, today's inflation, repair value, and other forward-looking objects are partly chosen to
honor constraints that private agents expected yesterday.

The full commitment residual vector on branch $b$ is
\begin{equation}
    \mathcal R_t^{C,b}
    =
    \left(
        \mathcal H_t^b,
        \mathcal R_t^{C,a},
        \mathcal R_t^{C,p}
    \right),
    \label{eq:app_deqn_commitment_full_residual}
\end{equation}
together with the branch-admissibility inequalities in \eqref{eq:app_opt_branch_admissibility}.
Compared with the discretionary vector \eqref{eq:app_deqn_discretion_full_residual}, the Bellman
value residual is replaced by Ramsey stationarity and promise-state recursion. Compared with the
rule-based vector \eqref{eq:app_deqn_rule_residual_vector}, the Taylor rule is replaced by
implementable optimality.

\subsection{Simulation and Moment Objects}

Once $F_\nu^P$ is trained, simulated paths are generated from
\begin{equation}
    \mathbf y_t^P=F_\nu^P(\mathcal Z_t^P),
    \qquad
    \mathcal Z_{t+1}^P=T^P(\mathcal Z_t^P,\varepsilon_{t+1}).
    \label{eq:app_deqn_simulation_law}
\end{equation}
The quantitative chapter can then compute the policy-relevant objects introduced in the main text:
\begin{equation}
    \mathcal O_t^P
    =
    \left(
        \Pi_t^P,
        \frac{Y_t^P}{Y_t^n(\mathcal Z_t^P)},
        \chi_t^P,
        I_t^{A,P},
        A_{t+1}^P
    \right),
    \label{eq:app_deqn_outcome_vector}
\end{equation}
and household welfare
\begin{equation}
    W^P(\mathcal Z_0^P)
    =
    \E_0
    \sum_{t=0}^{\infty}
    \beta^tU(C_t^P,N_t^P).
    \label{eq:app_deqn_welfare}
\end{equation}
These objects are not additional equilibrium conditions. They are the quantities used to report
stochastic steady states, ergodic distributions, impulse responses, and welfare comparisons across
policy environments.

\section{Expanded Optimal-Policy First-Order Conditions}
\label{app:expanded_optimal_policy_focs}

This appendix expands the key optimal-policy margins that are kept in compact form in
equations~\eqref{eq:discretion_foc_generic}--\eqref{eq:commitment_foc_generic}. The purpose is not to
replace the implementability system with a different model. It is to show explicitly how the same
private-equilibrium conditions generate the monetary-policy trade-off. Because the imported-input
cap and the non-negativity constraint on repair are complementarity restrictions, the differentiable
first-order conditions must be read branch by branch. The interior formulas are useful for seeing
the main margins; the Kuhn--Tucker branches below show what changes when the cap is slack or repair
is at the corner.

\subsection{Named Residuals}

For reference, define the residuals that matter for the policy margins:
\begin{align}
    G_t^E
    &\equiv
    1-\beta R_t
    \E_t
    \left[
        \frac{\Lambda_{t+1}}{\Lambda_t}
        \frac{1}{\Pi_{t+1}}
    \right],
    \label{eq:app_GE}\\
    G_t^Y
    &\equiv
    Y_t-C_t-p_t^mM_t-p_t^dS_t-p_{A,t}\Psi(I_t^A),
    \label{eq:app_GY}\\
    G_t^{A,I}
    &\equiv
    \Omega_t^Ap_{A,t}\Psi'(I_t^A)-Q_t^A,
    \label{eq:app_GAI}\\
    G_t^{A,L}
    &\equiv
    A_{t+1}-(1-\delta_A)A_t-I_t^A,
    \label{eq:app_GAL}\\
    G_t^Q
    &\equiv
    Q_t^A
    -
    \E_tM_{t,t+1}
    \left[
        -\C_{A,t+1}^{\Delta}
        +(1-\delta_A)Q_{t+1}^A
    \right],
    \label{eq:app_GQ}\\
    G_t^P
    &\equiv
    1-(1-\theta)(p_t^*)^{1-\varepsilon}
    -\theta\Pi_t^{\varepsilon-1},
    \label{eq:app_GP}\\
    G_t^\Delta
    &\equiv
    \Delta_t-(1-\theta)(p_t^*)^{-\varepsilon}
    -\theta\Pi_t^{\varepsilon}\Delta_{t-1}.
    \label{eq:app_GDelta}
\end{align}
The pricing-recursion residuals are
\begin{align}
    G_t^S
    &\equiv
    \mathcal S_t^p
    -
    \mc_tY_t
    -
    \theta\E_tM_{t,t+1}\Pi_{t+1}^{\varepsilon}\mathcal S_{t+1}^p,
    \label{eq:app_GS}\\
    G_t^F
    &\equiv
    \mathcal F_t^p
    -
    Y_t
    -
    \theta\E_tM_{t,t+1}\Pi_{t+1}^{\varepsilon-1}\mathcal F_{t+1}^p,
    \label{eq:app_GF}\\
    G_t^*
    &\equiv
    p_t^*
    -
    \frac{\varepsilon}{\varepsilon-1}
    \frac{\mathcal S_t^p}{\mathcal F_t^p}.
    \label{eq:app_Gstar}
\end{align}
Let $\mu_t^{D,j}$ and $\mu_t^{C,j}$ denote the current-value multipliers on residual $G_t^j$ under
discretion and commitment, respectively. The signs below depend only on the residual definitions
in \eqref{eq:app_GE}--\eqref{eq:app_Gstar}; changing the sign convention for residuals changes the
sign convention for the corresponding multipliers, not the economic content.

\subsection{Complementarity Branches}

The nonsmooth part of the baseline equilibrium can be written as two MCP restrictions:
\begin{equation}
    0\le \chi_t
    \perp
    K_t^M\equiv \bar m_t-M_t\ge0,
    \qquad
    0\le I_t^A
    \perp
    K_t^A\equiv \Omega_t^Ap_{A,t}\Psi'(I_t^A)-Q_t^A\ge0.
    \label{eq:app_mcp}
\end{equation}
The first condition says that the scarcity rent is positive only when the imported-input cap binds.
The second condition says that repair is positive only when the value of installed adaptation is high
enough to cover the private marginal installation cost. A local differentiable policy problem is
obtained by fixing the active branch:
\begin{align}
    \text{binding imported-input cap:}
    \qquad
    &G_t^M\equiv \bar m_t-M_t=0,
    \qquad
    \chi_t>0,
    \label{eq:app_binding_cap_branch}\\
    \text{slack imported-input cap:}
    \qquad
    &G_t^\chi\equiv \chi_t=0,
    \qquad
    \bar m_t-M_t>0,
    \label{eq:app_slack_cap_branch}\\
    \text{interior repair:}
    \qquad
    &G_t^{A,I}=0,
    \qquad
    I_t^A>0,
    \label{eq:app_interior_repair_branch}\\
    \text{repair corner:}
    \qquad
    &G_t^{I,0}\equiv I_t^A=0,
    \qquad
    G_t^{A,I}\ge0.
    \label{eq:app_corner_repair_branch}
\end{align}
Let $\mu_t^{P,M}$, $\mu_t^{P,\chi}$, and $\mu_t^{P,I,0}$ denote the active-branch multipliers in a
policy environment $P\in\{D,C\}$. On the binding-cap branch, the Lagrangian contains
$\mu_t^{P,M}G_t^M$. On the slack-cap branch, the cap equation is not imposed as an equality; instead
$\chi_t=0$ is imposed with multiplier $\mu_t^{P,\chi}$. On the interior-repair branch, the private
repair Euler condition $G_t^{A,I}=0$ is imposed. On the repair-corner branch, $I_t^A=0$ is imposed
and the private Euler inequality $G_t^{A,I}\ge0$ is slack unless the economy is exactly at the
switching boundary.

The two imported-input branches have different monetary-policy content. If the cap binds, aggregate
desired demand satisfies
\begin{equation}
    M_t^d(Y_t,p_t^m+\chi_t,A_t)=\bar m_t.
    \label{eq:app_binding_import_demand}
\end{equation}
Holding the external state and installed adaptation fixed, differentiation with respect to the
policy rate gives
\begin{equation}
    M_Y\frac{\partial Y_t}{\partial R_t}
    +
    M_p\frac{\partial \chi_t}{\partial R_t}
    =
    0,
    \qquad
    M_Y>0,\quad M_p<0.
    \label{eq:app_binding_chi_derivative}
\end{equation}
Thus the current bottleneck-relief result is a binding-cap result: if a contractionary policy
innovation lowers current output, it lowers desired use of the constrained imported input and hence
reduces the scarcity rent. If the cap is slack, the active equation is instead
\begin{equation}
    \chi_t=0,
    \qquad
    \frac{\partial\chi_t}{\partial R_t}=0
    \quad
    \text{locally as long as the slack branch is not left.}
    \label{eq:app_slack_chi_derivative}
\end{equation}
In that region monetary policy can still affect inflation through ordinary demand and cost channels,
but there is no current scarcity-rent channel because scarce access has zero equilibrium price.

The repair branches change the direct policy margin. On the interior branch, the direct $R_t$
condition contains the repair-cost term shown below in \eqref{eq:app_discretion_R} and
\eqref{eq:app_commitment_R}. On a strict repair-corner branch, $G_t^{A,I}>0$ and the private repair
Euler inequality is not an equality. The active repair equation is $I_t^A=0$, so the direct
installation-cost term drops out of the $R_t$ first-order condition:
\begin{equation}
    -\mu_t^{P,E}
    \beta
    \E_t
    \left[
        \frac{\Lambda_{t+1}}{\Lambda_t}
        \frac{1}{\Pi_{t+1}}
    \right]
    =
    0,
    \qquad
    P\in\{D,C\},
    \label{eq:app_corner_R_condition}
\end{equation}
provided the economy remains strictly inside the no-repair region. This does not mean that monetary
policy is irrelevant for repair forever. It means that, locally at a strict corner, a small policy
change does not change current repair because the firm is not marginally investing. Policy can still
move the economy toward or away from the switching boundary through $Q_t^A$ and future states.

The current repair-investment first-order condition is also branch dependent. The interior branch
uses \eqref{eq:app_discretion_I} or \eqref{eq:app_commitment_I}. At a strict corner, replace the
private repair Euler equality by $G_t^{I,0}=I_t^A=0$. The active-branch condition is then
\begin{equation}
    -\mu_t^{P,Y}p_{A,t}\Psi'(0)
    -
    \mu_t^{P,A,L}
    +
    \mu_t^{P,I,0}
    =
    0,
    \qquad
    G_t^{A,I}\big|_{I_t^A=0}
    =
    \Omega_t^Ap_{A,t}-Q_t^A
    \ge0.
    \label{eq:app_corner_I_condition}
\end{equation}
At the switching boundary, $I_t^A=0$ and $G_t^{A,I}=0$ hold simultaneously. The correct object is
then the mixed-complementarity system \eqref{eq:app_mcp}, not a single smooth derivative. This is why
the local sign result for repair in Proposition~\ref{prop:baseline_tradeoff} is stated for the
interior branch, while Corollary~\ref{cor:no_repair} describes the no-repair branch.
This does not contradict the repair trade-off in the baseline. The trade-off is active when firms
are marginally choosing how much to invest in reducing imported-input dependence. At a strict
no-repair corner, current repair is not marginal: the firm is already choosing zero because the
private value of installed adaptation is below its marginal cost. In that region, monetary policy
affects repair only by changing the distance to the switching boundary, through the continuation
value $Q_t^A$ and the future bottleneck state. Once the inequality becomes tight, the economy moves
back to the interior or boundary case where the repair margin again enters the policy FOCs.

\subsection{Discretion}

Under discretion, the central bank chooses a current implementable allocation taking future
discretionary policy functions as given. The implementation rate has no direct utility payoff, so
its direct first-order condition is obtained from the two implementability constraints in which
$R_t$ enters directly:
\begin{equation}
    -\mu_t^{D,E}
    \beta
    \E_t
    \left[
        \frac{\Lambda_{t+1}}{\Lambda_t}
        \frac{1}{\Pi_{t+1}}
    \right]
    +
    \mu_t^{D,A,I}
    \vartheta_Ap_{A,t}\Psi'(I_t^A)
    =
    0.
    \label{eq:app_discretion_R}
\end{equation}
The first term is the value of using $R_t$ to implement intertemporal demand. The second term is
the value of the same instrument through the private cost of repair. If $\vartheta_A=0$, the direct
repair-cost term disappears and policy affects repair only through equilibrium continuation values.

The first-order condition for current adaptation investment on the interior branch is
\begin{equation}
    -\mu_t^{D,Y}p_{A,t}\Psi'(I_t^A)
    +
    \mu_t^{D,A,I}\Omega_t^Ap_{A,t}\Psi''(I_t^A)
    -
    \mu_t^{D,A,L}
    =
    0.
    \label{eq:app_discretion_I}
\end{equation}
This expression separates the three roles of $I_t^A$. It absorbs current real resources through
$G_t^Y$, changes the private repair Euler equation through $G_t^{A,I}$, and moves the next
adaptation stock through $G_t^{A,L}$. This is the formal version of the statement that repair is
both a current expenditure and a change in future production capability.

The first-order condition for installed next-period adaptation is
\begin{equation}
    \beta
    \E_t
    \left[
        V_A^D(\mathcal Z_{t+1})
    \right]
    +
    \mu_t^{D,A,L}
    +
    \mu_t^{D,Q}
    \frac{\partial G_t^Q}{\partial A_{t+1}}
    =
    0.
    \label{eq:app_discretion_A_next}
\end{equation}
The first term is the policymaker's value of entering the next period with more installed
adaptation. The second term is the law-of-motion multiplier. The third term appears because
private firms value adaptation through the continuation-value recursion; changing $A_{t+1}$ changes
the implementability condition that defines $Q_t^A$.

For current inflation, the direct discretionary condition is pinned down by the Calvo price-index
and dispersion equations:
\begin{equation}
    -\mu_t^{D,P}\theta(\varepsilon-1)\Pi_t^{\varepsilon-2}
    -
    \mu_t^{D,\Delta}\theta\varepsilon\Pi_t^{\varepsilon-1}\Delta_{t-1}
    =
    0.
    \label{eq:app_discretion_pi}
\end{equation}
Inflation therefore does not enter the policy problem as a primitive quadratic loss. It enters
because it must be consistent with the Calvo price index and because it changes the next price
dispersion state, which raises gross production requirements.

Finally, the first-order condition for price dispersion itself is
\begin{equation}
    \beta
    \E_t
    \left[
        V_\Delta^D(\mathcal Z_{t+1})
    \right]
    +
    \mu_t^{D,\Delta}
    +
    \mathfrak D_t^{D,\Delta}
    =
    0.
    \label{eq:app_discretion_Delta}
\end{equation}
Here $\mathfrak D_t^{D,\Delta}$ collects the current implementability derivatives of the
cost-minimizing input-demand and resource constraints with respect to $\Delta_t$. This term is not
zero because aggregate input requirements are proportional to $\Delta_tY_t$. Thus the policy value
of inflation stabilization is connected to real production costs through price dispersion.

\subsection{Commitment}

Under commitment, the current direct margins remain present, but forward-looking private
constraints add promise terms. With the current-value normalization used in
\eqref{eq:commitment_foc_generic}, the commitment analogue of
\eqref{eq:app_discretion_R} is
\begin{equation}
    -\mu_t^{C,E}
    \beta
    \E_t
    \left[
        \frac{\Lambda_{t+1}}{\Lambda_t}
        \frac{1}{\Pi_{t+1}}
    \right]
    +
    \mu_t^{C,A,I}
    \vartheta_Ap_{A,t}\Psi'(I_t^A)
    =
    0,
    \label{eq:app_commitment_R}
\end{equation}
unless the chosen recursive implementation gives $R_t$ an additional role in future constraints.
The deeper difference between discretion and commitment appears for variables that enter
forward-looking conditions or become future states.

For example, current inflation enters the previous-period household Euler equation and Calvo
pricing recursions as a future object. The commitment inflation condition is therefore
\begin{equation}
    -\mu_t^{C,P}\theta(\varepsilon-1)\Pi_t^{\varepsilon-2}
    -
    \mu_t^{C,\Delta}\theta\varepsilon\Pi_t^{\varepsilon-1}\Delta_{t-1}
    +
    \varrho_{t-1}^{C,\Pi}
    =
    0,
    \label{eq:app_commitment_pi}
\end{equation}
where the lagged promise term is
\begin{align}
    \varrho_{t-1}^{C,\Pi}
    &=
    p_t^{C,E}
    \frac{\partial \widetilde G_{t-1}^{E}}{\partial \Pi_t}
    +
    p_t^{C,S}
    \frac{\partial \widetilde G_{t-1}^{S}}{\partial \Pi_t}
    +
    p_t^{C,F}
    \frac{\partial \widetilde G_{t-1}^{F}}{\partial \Pi_t}.
    \label{eq:app_pi_promise_term}
\end{align}
The tildes denote the scaled forward-looking constraints. The previous-period marginal-utility
normalizer is included in the promise variables $p_t^{C,E}$, $p_t^{C,S}$, and $p_t^{C,F}$, so the
current derivatives depend on the current marginal utility $\Lambda_t=C_t^{-\sigma}$ but do not
require a separate lagged-consumption state. Under the normalization used in the computational
system, the relevant current-price derivatives are
\begin{align}
    \frac{\partial \widetilde G_{t-1}^{E}}{\partial \Pi_t}
    &=
    -\Lambda_t\Pi_t^{-2},
    \label{eq:app_GE_pi_lag}\\
    \frac{\partial \widetilde G_{t-1}^{S}}{\partial \Pi_t}
    &=
    \Lambda_t\varepsilon\Pi_t^{\varepsilon-1}\mathcal S_t^p,
    \label{eq:app_GS_pi_lag}\\
    \frac{\partial \widetilde G_{t-1}^{F}}{\partial \Pi_t}
    &=
    \Lambda_t(\varepsilon-1)\Pi_t^{\varepsilon-2}\mathcal F_t^p.
    \label{eq:app_GF_pi_lag}
\end{align}
This is the formal source of the expectations channel under commitment: current inflation is priced
not only through current price dispersion, but also through promises embedded in earlier
forward-looking constraints.

The commitment condition for current repair investment is
\begin{equation}
    -\mu_t^{C,Y}p_{A,t}\Psi'(I_t^A)
    +
    \mu_t^{C,A,I}\Omega_t^Ap_{A,t}\Psi''(I_t^A)
    -
    \mu_t^{C,A,L}
    =
    0,
    \label{eq:app_commitment_I}
\end{equation}
while the condition for installed next-period adaptation is
\begin{equation}
    \mu_t^{C,A,L}
    +
    \mu_t^{C,Q}
    \frac{\partial G_t^Q}{\partial A_{t+1}}
    +
    \E_t
    \left[
        \beta
        (\boldsymbol\mu_{t+1}^{C})'
        \frac{\partial \mathcal G_{t+1}^{B}}{\partial A_{t+1}}
    \right]
    =
    0.
    \label{eq:app_commitment_A_next}
\end{equation}
The last term is absent from the direct discretionary repair condition. It is the Ramsey value of
how installed adaptation changes all future implementability constraints: future marginal costs,
future scarcity rents, future price dispersion, and future private repair values. This is why
commitment can change current repair incentives even when the contemporaneous private installation
technology is unchanged.

\end{document}
